% LuaLaTeX document — compile with: lualatex cube_galerkin27.tex
\documentclass[11pt,a4paper]{article}

\usepackage{fontspec}
\setmainfont{Latin Modern Roman}
\setmonofont{Latin Modern Mono}

\usepackage{amsmath,amssymb,amsthm,bm}
\usepackage{booktabs}
\usepackage{geometry}
\geometry{margin=2.5cm}
\usepackage{hyperref}
\hypersetup{colorlinks=true,linkcolor=blue!60!black,citecolor=green!50!black,urlcolor=blue!70!black}
\usepackage{enumitem}
\usepackage{xcolor}
\usepackage{mathtools}
\usepackage{graphicx}  % \resizebox for long closed-form expressions

\newcommand{\D}{\Delta}
\newcommand{\order}{\mathcal{O}}
\newcommand{\bG}{\bm{G}}
\newcommand{\bs}{\bm{s}}
\newcommand{\br}{\bm{r}}
\newcommand{\bxi}{\bm{\xi}}
\newcommand{\be}{\bm{\varepsilon}}
\newcommand{\bu}{\bm{u}}
\newcommand{\bphi}{\bm{\varphi}}
\newcommand{\R}{\mathbb{R}}
\newcommand{\PV}{\mathrm{p.v.}}
\newcommand{\indV}{\chi_V}
\newcommand{\Lspace}{L^1_{\mathrm{loc}}}

\theoremstyle{plain}
\newtheorem{theorem}{Theorem}
\newtheorem{proposition}[theorem]{Proposition}
\newtheorem{lemma}[theorem]{Lemma}
\theoremstyle{definition}
\newtheorem{definition}[theorem]{Definition}
\theoremstyle{remark}
\newtheorem{remark}[theorem]{Remark}

\title{A Rigorous Galerkin Closure for the Elastodynamic\\
Volume Integral Equation on a Cube\\
\large 27-Component Path-B Formulation}
\author{Research Notes}
\date{April 2026}

\begin{document}
\maketitle
\tableofcontents
\vspace{1em}

%======================================================================
\section{Motivation: differentiating under the integral fails}
%======================================================================

{\sloppy
The classical moment-method closure for the cube T-matrix begins from the
Lippmann--Schwinger (LS) equation\par}
\begin{equation}
  u_i(\br) \;=\; u_i^{0}(\br)
   + \omega^2 \!\int_V G_{ij}(\br - \br')\,\delta\rho(\br')\, u_j(\br')\,dV'
   - \!\int_V \partial'_k G_{in}(\br - \br')\,\delta c_{nklj}(\br')\,
   \partial'_l u_j(\br')\,dV',
   \label{eq:LS}
\end{equation}
where $G_{ij}$ is the elastostatic Green's tensor and $V \subset \R^3$ is the
inclusion. To obtain a closed finite-dimensional system one then formally
takes derivatives at the centre $\br = \mathbf{0}$:
\begin{align}
   u_i(\mathbf{0}) &= u_i^{0}(\mathbf{0}) + \omega^2\!\int_V G_{ij} \,\delta\rho\, u_j\, dV'
                       - \!\int_V \partial'_k G_{in}\,\delta c_{nklj}\,\partial'_l u_j\, dV', \nonumber\\
   \partial_p u_i(\mathbf{0}) &= \partial_p u_i^{0}(\mathbf{0})
                       - \omega^2\!\int_V \partial'_p G_{ij} \,\delta\rho\, u_j\, dV'
                       + \!\int_V \partial'_p \partial'_k G_{in}\,\delta c_{nklj}\,\partial'_l u_j\, dV', \label{eq:LS-derivs}\\
   \partial_q\partial_p u_i(\mathbf{0}) &= \partial_q\partial_p u_i^{0}(\mathbf{0})
                       + \omega^2\!\int_V \partial'_q\partial'_p G_{ij} \,\delta\rho\, u_j\, dV'
                       - \!\int_V \partial'_q\partial'_p\partial'_k G_{in}\,\delta c_{nklj}\,\partial'_l u_j\, dV',
   \nonumber
\end{align}
where the $r$--derivatives have been converted to $r'$--derivatives by translation
invariance $\partial_r G(\br - \br') = -\partial_{r'} G(\br - \br')$.

\medskip

The Leibniz rule used in passing from \eqref{eq:LS} to \eqref{eq:LS-derivs}
requires (DCT version) a dominating function $|\partial_r f(\br,\br')| \le g(\br')$
with $g \in L^{1}(V)$.  The static elastic kernel has the singularity
\begin{equation}
   G^{\text{stat}}_{ij}(\bs) \sim \frac{1}{|\bs|},
   \quad \partial' G \sim \frac{1}{|\bs|^2},
   \quad \partial'\partial' G \sim \frac{1}{|\bs|^3},
   \quad \partial'\partial'\partial' G \sim \frac{1}{|\bs|^4},
   \label{eq:singularities}
\end{equation}
so already at the level of $\partial' G$ \emph{no integrable dominating function exists} when $\br$ is in the support of the inclusion. Each of the equations in
\eqref{eq:LS-derivs} therefore makes only \emph{formal} sense; converting them into rigorous statements requires distributional regularization (Cauchy principal value, Hadamard finite part, and Eshelby-type delta corrections).

\begin{table}[h]
\centering
\begin{tabular}{lll}
\toprule
LS derivative order & worst kernel & status \\
\midrule
$0$ (body force)        & $1/|\bs|$   & $\Lspace$ \checkmark \\
$0$ (elastic kernel)    & $1/|\bs|^2$ & conditionally convergent (PV) \\
$1$ (body force)        & $1/|\bs|^2$ & PV \\
$1$ (elastic kernel)    & $1/|\bs|^3$ & \textbf{Eshelby delta required} \\
$2$ (body force)        & $1/|\bs|^3$ & \textbf{Eshelby delta required} \\
$2$ (elastic kernel)    & $1/|\bs|^4$ & \textbf{derivative-of-delta / Hadamard PF} \\
\bottomrule
\end{tabular}
\caption{Singular structure of the formal closure equations \eqref{eq:LS-derivs}.}
\label{tab:singularities}
\end{table}

%======================================================================
\section{Path B: Galerkin formulation on a polynomial trial space}
%======================================================================

We replace point evaluation by a Galerkin (weak) projection onto a
finite-dimensional polynomial trial space.  The resulting bilinear forms
involve only $G_{ij}$ itself (the singularity $1/|\bs|$ is locally
$L^{1}$ in $\R^{3}$), so \emph{every entry of the closure matrix is an
absolutely convergent integral}.  No principal values, no delta
corrections, no distributional bookkeeping.

\subsection{Trial space \texorpdfstring{$T_{27}$}{T27}}

Let $V = [-a,a]^{3}$.  Denote by $P_{2}(V)$ the space of real polynomials
in $\br = (r_1,r_2,r_3)$ of total degree $\le 2$, $\dim P_{2}(V) = 10$. The naive
vector trial space $(P_{2}(V))^{3}$ has dimension $30$.  We remove the
$3$ rigid-rotation modes by symmetrizing the linear part, giving the
$27$-dimensional subspace
\begin{equation}
   T_{27} \;=\; \Big\{\,\bphi(\br)\;\Big|\;
   \varphi_i(\br) = U_i + \varepsilon_{ai}\, r_a + \tfrac{1}{2}\,M_{abi}\, r_a r_b,\quad
   \varepsilon_{ai} = \varepsilon_{ia},\quad M_{abi} = M_{bai}\,\Big\}.
\end{equation}
The dimensions add as $3 + 6 + 18 = 27$, matching one displacement
monopole, the symmetric strain tensor, and the symmetric second-derivative
tensor of $\bu$ (equivalent to the strain gradient $\partial_k\varepsilon_{ij}$
through the kinematic identity
$\partial_k\partial_l u_m = \partial_k\varepsilon_{lm} + \partial_l\varepsilon_{km} - \partial_m\varepsilon_{kl}$).

\subsection{Galerkin LS}

\begin{definition}[Path-B closure]
The Path-B 27-component closure is the variational problem:
find $\bu^{h}\in T_{27}$ such that for all $\bphi \in T_{27}$,
\begin{equation}
   \langle\bphi,\bu^{h}\rangle_V
   \;=\;
   \langle\bphi,\bu^{0}\rangle_V
   + \omega^{2}\delta\rho\, \mathcal{B}_{\text{body}}[\bphi,\bu^{h}]
   - \delta c_{nklj}\, \mathcal{B}_{\text{el}}[\bphi,\bu^{h}],
   \label{eq:galerkin-LS}
\end{equation}
where the bilinear forms are
\begin{align}
\langle\bphi,\bu\rangle_V &:= \int_{V} \varphi_i(\br)\, u_i(\br)\,dV, \\[2pt]
\mathcal{B}_{\text{body}}[\bphi,\bu] &:= \int_{V}\!\!\int_{V} \varphi_i(\br)\,
            G_{in}(\br-\br')\, u_n(\br')\,dV\,dV', \\[2pt]
\mathcal{B}_{\text{el}}[\bphi,\bu] &:= \int_{V}\!\!\int_{V} \varphi_i(\br)\,
            \partial'_k G_{in}(\br-\br')\, \partial'_l u_j(\br')\,dV\,dV'.
\end{align}
\end{definition}

\subsection{Why this is well defined}

\begin{proposition}\label{prop:wellposed}
For polynomial $\bphi,\bu \in T_{27}$, every bilinear form in
\eqref{eq:galerkin-LS} is an absolutely convergent integral.  No
distributional regularization is needed.
\end{proposition}
\begin{proof}[Proof sketch]
The mass form is trivial because the integrand is a polynomial.
For $\mathcal{B}_{\text{el}}$, we use integration by parts on $\bu$
(legitimate because $\bu$ is a polynomial and the boundary term is a
2D integral against a smooth $L^1_{\text{loc}}$ kernel):
\begin{equation}
\mathcal{B}_{\text{el}}[\bphi,\bu]
=
-\int_{V}\!\!\int_V\!\varphi_i\,G_{in}\,\partial'_k\partial'_l u_j\,dV'\,dV
+ \int_{V}\!\!\oint_{\partial V}\!\varphi_i\,G_{in}\,\partial'_l u_j\,n_k\,dS'\,dV.
\end{equation}
{\sloppy
Both integrands now contain only $G_{in}(\br-\br') \sim 1/|\br-\br'|$ multiplied
by polynomials/smooth functions.  The 6-fold integral
$\int\!\!\int_{V\times V}1/|\br-\br'|\,dV\,dV'$ converges by the
Hardy--Littlewood--Sobolev inequality (or directly: in
$\R^{6}$ the singular set $\{\br = \br'\}$ has codimension $3$ and
$1/|\br-\br'|$ has order-1 singularity, hence is locally
$L^{1}$ in $\R^{6}$).  The 5-fold boundary integral is similar.
The body form is even simpler since no integration by parts is needed.\par}
\qed
\end{proof}

\subsection{Equivalence to point-collocation in the limit}

The classical Path-A closure (point evaluation at $\br = \mathbf{0}$) is
the formal limit of a Petrov--Galerkin scheme with Dirac-delta test
functions.  Both Path A and Path B converge to the exact volume integral
equation as the polynomial degree $\to \infty$; at finite truncation
their stiffness matrices differ but encode the same physical operator.
Path B is preferable for analysis because it is symmetric, energy
conserving, and free of singular regularizations.

%======================================================================
\section{Center-of-mass / relative coordinates}
\label{sec:cm-relative}
%======================================================================

The 6-fold integrals in \eqref{eq:galerkin-LS} are reduced to 3-fold
integrals against the \emph{cube autocorrelation} via the \emph{centre-of-mass and relative coordinate} change of variables (also called \emph{sum-and-difference variables} in the Eshelby/Mura literature, or \emph{Jacobi coordinates} in the simplest 2-body case):
\begin{equation}
   \bxi = \tfrac{1}{2}(\br + \br'), \qquad \bs = \br - \br',
   \qquad \br = \bxi + \tfrac{\bs}{2}, \quad \br' = \bxi - \tfrac{\bs}{2},
   \qquad d\br\,d\br' = d\bxi\,d\bs.
\end{equation}
The constraint $\br,\br' \in V$ becomes $\bxi \in V \cap (V + \bs/2)\cap (V - \bs/2)$, which for the cube $V = [-a,a]^{3}$ simplifies to a centred box of half-lengths
$b_i(\bs) = a - |s_i|/2$:
\begin{equation}
   \bxi \in B(\bs) := \prod_{i=1}^{3} \big[-(a - |s_i|/2),\, a - |s_i|/2\big],
   \qquad \text{valid when } |s_i| < 2a.
\end{equation}
The 6-fold integral becomes
\begin{equation}
   \int_{V}\!\!\int_{V} P(\br)\,K(\br - \br')\,Q(\br')\,d\br\,d\br'
   \;=\; \int_{[-2a,2a]^{3}}\! K(\bs)
        \underbrace{\Big[\int_{B(\bs)}\!P\!\big(\bxi+\tfrac{\bs}{2}\big)\,Q\!\big(\bxi-\tfrac{\bs}{2}\big)\,d\bxi\Big]}_{\text{polynomial in } (\bs, a)}\,d\bs.
   \label{eq:cm-relative}
\end{equation}
For polynomial $P,Q$ the inner integral is elementary: every $\xi_i^{p}$
integrates over $[-b_i, b_i]$ to $\frac{2 b_i^{p+1}}{p+1}$ when $p$ is even and to $0$
when $p$ is odd, giving a polynomial in
$(b_1,b_2,b_3) = (a-|s_1|/2,\,a-|s_2|/2,\,a-|s_3|/2)$.

\medskip

\noindent\textbf{Cube autocorrelation.} For $P = Q = 1$ the inner integral
is exactly the autocorrelation of the cube indicator,
\begin{equation}
   (\indV \star \indV)(\bs) \;=\; \prod_{i=1}^{3} \big(2a - |s_i|\big)_{+},
\end{equation}
the 3-D \emph{tent function} $\Lambda_V(\bs)$.

%======================================================================
\section{Master integrals}
%======================================================================

After the centre-of-mass/relative reduction, every entry of the
$27\times 27$ Galerkin matrix is a finite linear combination of the
following \emph{master integrals} over the doubled cube
$\Omega = [-2a,2a]^{3}$:
\begin{align}
   \mathcal{I}_{p,q,r}(a) &:= \int_{\Omega}\!\frac{s_1^{p}\,s_2^{q}\,s_3^{r}}{|\bs|}\,d\bs, \label{eq:Imaster}\\
   \mathcal{J}^{(ij)}_{p,q,r}(a) &:= \int_{\Omega}\!\frac{s_i s_j\,s_1^{p}\,s_2^{q}\,s_3^{r}}{|\bs|^{3}}\,d\bs.
   \label{eq:Jmaster}
\end{align}
For the $9$-component (current) closure only $p+q+r \le 3$ is needed; the
$27$-component closure requires $p+q+r \le 5$.

By the parity $\bs \to -\bs$ on each axis separately, $\mathcal{I}_{p,q,r}$
vanishes whenever any of $p,q,r$ is odd, and $\mathcal{J}^{(ij)}_{p,q,r}$
vanishes unless the multiset $\{i,j,p\text{-copies of }1,q\text{-copies of }2,r\text{-copies of }3\}$ has even count for each axis.  This collapses
the formally $\binom{8}{3} \cdot 3 = 168$ entries down to $\sim 20$ inequivalent integrals after cubic symmetry.

\subsection{Tier-1 results (computed analytically)}\label{sec:tier1}

The basic integral $\mathcal{I}_{0,0,0}(a)$ is the dimensionless cube
constant $g_0$ scaled by $a^{2}$.  Mathematica's \texttt{Integrate}
returns it in the form
\begin{equation}
   \int_{[-1,1]^{3}}\!\frac{1}{|\bu|}\,du \;=\; -2\Big(\pi + \log\!\big(1351 - 780\sqrt{3}\big)\Big),
\end{equation}
which equals the Pell-stable form used in the existing
\texttt{CubeAnalytic.wl} package via the identity
$1351 - 780\sqrt{3} = (70226 + 40545\sqrt{3})^{-2/3}$ (both sides are Pell
units of $\mathbb{Z}[\sqrt{3}]$, satisfying $1351^{2} - 3\cdot 780^{2} = 1$
and $70226^{2} - 3\cdot 40545^{2} = 1$). Numerical agreement to $\sim\!78$ digits has been verified.

% Auto-generated by Mathematica/CubeGalerkin27.wl with hand-edited
% \begin{aligned}[t]...\end{aligned} line breaks in Tables 2--3 to prevent
% long closed-form Mp expressions from overflowing the page. If you
% regenerate this file, reapply the aligned[t] wrappers to the overflowing
% rows (see git log for this file).
\subsection{Verified analytical master integrals}
\label{sec:results}

The verification at degree 0 was reported in Section~\ref{sec:tier1}.
We extend the table to include all parity-allowed master integrals up to degree~6:

\begin{table}[h]
\centering
\begin{tabular}{lll}
\toprule
$(p,q,r)$ & $\mathcal{M}_{p,q,r} = \int_{[-1,1]^3}\!\frac{x^p y^q z^r}{|\bm{r}|}\,dV$ & numerical \\
\midrule
$(0,0,0)$ & $-2\pi-\tfrac{4}{3}\log(70226-40545\sqrt{3})$ & $9.5203095$ \\
$(2,0,0)$ & $\tfrac{1}{9}\bigl(6\sqrt{3}-\pi+\log(3650401+2107560\sqrt{3})\bigr)$ & $2.5615786$ \\
$(2,2,0)$ & $\begin{aligned}[t]&\tfrac{1}{135}\bigl(108\sqrt{3}+18\pi+63\log 2+20\log(26-15\sqrt{3})\\&\quad{}-60\log(\sqrt{3}-1)-66\log(1+\sqrt{3})-90\coth^{-1}\sqrt{3}\bigr)\end{aligned}$ & $0.75095335$ \\
$(4,0,0)$ & $\begin{aligned}[t]&-\tfrac{2}{45}\bigl(12\sqrt{3}+7\pi-42\log 2\\&\quad{}+99\log(\sqrt{3}-1)-15\log(1+\sqrt{3})\bigr)\end{aligned}$ & $1.4351485$ \\
$(2,2,2)$ & $\begin{aligned}[t]&\tfrac{1}{1260}\bigl(1008\sqrt{3}-42\pi+522\log 2+441\log(\sqrt{3}-1)\\&\quad{}-1395\log(1+\sqrt{3})-84\log(2+\sqrt{3})-90\log(11\sqrt{3}-19)\\&\quad{}+35\log(26+15\sqrt{3})-672\coth^{-1}\sqrt{3}\bigr)\end{aligned}$ & $0.22735457$ \\
$(4,2,0)$ & $\begin{aligned}[t]&\tfrac{1}{630}\bigl(-48\sqrt{3}+17\pi\\&\quad{}+12\tanh^{-1}\tfrac{21252634831\sqrt{3}}{36810643322}\bigr)\end{aligned}$ & $0.42942069$ \\
$(6,0,0)$ & $\tfrac{1}{42}\bigl(-5\pi-4\log 2+8(3\sqrt{3}+\log(5+3\sqrt{3}))\bigr)$ & $0.99201788$ \\
\bottomrule
\end{tabular}
\caption{Closed-form values of the scalar master integrals $\mathcal{M}_{p,q,r}$ over the unit cube.  Cubic symmetry equates entries that are permutations of each other; only canonical orderings $p\ge q \ge r$ are listed.}
\label{tab:tier1-masters}
\end{table}

The same family of integrals over the positive octant $[0,1]^3$ supplies the body-bilinear-form building blocks; we collect them in Table~\ref{tab:tier2-masters}.

\begin{table}[h]
\centering
\begin{tabular}{lll}
\toprule
$(p,q,r)$ & $\mathcal{M}^+_{p,q,r} = \int_{[0,1]^3}\!\frac{x^p y^q z^r}{|\bm{r}|}\,dV$ & numerical \\
\midrule
$(0,0,0)$ & $-\frac{\pi }{4}-\frac{1}{6} \log \left(70226-40545 \sqrt{3}\right)$ & $1.1900387$ \\
$(1,0,0)$ & $\begin{aligned}[t]&\tfrac{1}{36}\bigl(-2\pi-12\bigl(\sqrt{2}-\sqrt{3}+\sinh^{-1}(1)\bigr)\\&\quad{}+\cosh^{-1}(189750626)+18\coth^{-1}\sqrt{3}\bigr)\end{aligned}$ & $0.51559356$ \\
$(1,1,0)$ & $\begin{aligned}[t]&\tfrac{1}{12}\bigl(1-7\sqrt{2}+6\sqrt{3}+\log(5\sqrt{2}-7)\\&\quad{}-3\log(2-\sqrt{3})\bigr)\end{aligned}$ & $0.23329690$ \\
$(2,0,0)$ & $\begin{aligned}[t]&\tfrac{1}{144}\bigl(12\sqrt{3}-2\pi\\&\quad{}+\log(26650854921601+15386878263120\sqrt{3})\bigr)\end{aligned}$ & $0.32019732$ \\
$(1,1,1)$ & $\tfrac{1}{5}\bigl(1-4\sqrt{2}+3\sqrt{3}\bigr)$ & $0.10785963$ \\
$(2,1,0)$ & $\begin{aligned}[t]&\tfrac{1}{2880}\bigl(-504\sqrt{2}+672\sqrt{3}+32\pi+72\log(\sqrt{3}-1)\\&\quad{}-252\log(1+\sqrt{3})+90\log(2(2+\sqrt{3}))-261\sinh^{-1}(1)\\&\quad{}+40\cosh^{-1}(26)+45\coth^{-1}\sqrt{2}\bigr)\end{aligned}$ & $0.14741068$ \\
\bottomrule
\end{tabular}
\caption{Closed-form values of the positive-octant master integrals $\mathcal{M}^+_{p,q,r}$ over $[0,1]^3$. Up to canonical ordering $p\ge q \ge r$ and $p+q+r\le 3$ they exhaust all polynomial weights needed by the Path-B body bilinear form on the constant-and-linear subspace.}
\label{tab:tier2-masters}
\end{table}

\subsection{Body bilinear-form building blocks}
\label{sec:results-tier3}

For the constant subspace of $T_{27}$ the diagonal and off-diagonal entries of the body bilinear form $\mathcal{B}_{\mathrm{body}}$ reduce, after the centre-of-mass / relative-coordinate substitution and an octant-symmetry restriction, to three positive scalars.
Define
\begin{align}
  i_A      &:= \int_{[0,1]^3}\!(1-u_1)(1-u_2)(1-u_3)\,\frac{1}{|\bm{u}|}\,d\bm{u},\\
  i_{B,11} &:= \int_{[0,1]^3}\!(1-u_1)(1-u_2)(1-u_3)\,\frac{u_1^2}{|\bm{u}|^3}\,d\bm{u},\\
  i_{B,12} &:= \int_{[0,1]^3}\!(1-u_1)(1-u_2)(1-u_3)\,\frac{u_1 u_2}{|\bm{u}|^3}\,d\bm{u}.
\end{align}
The first reduces by linearity to a signed combination of $\mathcal{M}^+_{p,q,r}$ values (Table~\ref{tab:tier2-masters}):
\begin{equation}
  i_A \;=\; \mathcal{M}^+_{0,0,0} \;-\; 3\mathcal{M}^+_{1,0,0} \;+\; 3\mathcal{M}^+_{1,1,0} \;-\; \mathcal{M}^+_{1,1,1}
       \;\approx\; 0.2352890805487075.
  \label{eq:iA-closed}
\end{equation}

The remaining two are obtained by integrating once by parts in the $u_1$ direction; the singular kernel $|\bm{u}|^{-3}$ collapses to a $|\bm{u}|^{-1}$-class kernel and a clean two-dimensional residual:
\begin{align}
  i_{B,11} &\;=\; i_A \;-\; J_{11},\qquad
  J_{11} \;=\; \int_{[0,1]^2}\!(1-u_2)(1-u_3)\Bigl[\sqrt{1+u_2^2+u_3^2} - \sqrt{u_2^2+u_3^2}\,\Bigr]\,du_2\,du_3,\\
  i_{B,12} &\;=\; \int_{[0,1]^2}\!(1-u_2)(1-u_3)\,u_2\,\Bigl[\frac{1}{\sqrt{u_2^2+u_3^2}} \;-\; \operatorname{arcsinh}\!\Bigl(\frac{1}{\sqrt{u_2^2+u_3^2}}\Bigr)\Bigr]\,du_2\,du_3.
\end{align}

Mathematica evaluates both two-dimensional integrals in closed form:
\begin{equation*}
\resizebox{\linewidth}{!}{$\displaystyle
J_{11} \;=\; \frac{1}{240} \left(8+8 \sqrt{2}-16 \sqrt{3}-55 \log (2)+15 \log \left(3+2 \sqrt{2}\right)+110 \log \left(1+\sqrt{3}\right)-15 \log \left(2+\sqrt{3}\right)-80 \tan ^{-1}\left(\frac{3}{8}\right)+80 \tan ^{-1}\left(\frac{1}{183} \left(96-73 \sqrt{3}\right)\right)+10 \sinh ^{-1}(1)\right)
$}
\end{equation*}
with $J_{11} \approx 0.1568593870324717$, and
\begin{equation*}
\resizebox{\linewidth}{!}{$\displaystyle
i_{B,12} \;=\; \frac{1}{240} \left(-8 \sqrt{2}+4 \sqrt{3}+20 \log (2)+5 \log (256)+40 \log \left(1+\sqrt{2}\right)+22 \log \left(3+2 \sqrt{2}\right)-120 \log \left(1+\sqrt{3}\right)-22 \log \left(2+\sqrt{3}\right)+20 \tan ^{-1}\left(\frac{1}{2}\right)+16 \tan ^{-1}\left(\frac{5}{8}\right)-40 \tan ^{-1}\left(\frac{3}{2}\right)+40 \tan ^{-1}\left(8+\frac{13}{\sqrt{3}}\right)-20 \tan ^{-1}\left(\frac{1}{11} \left(8-5 \sqrt{3}\right)\right)-16 \tan ^{-1}\left(\frac{1}{167} \left(160+89 \sqrt{3}\right)\right)+4 \sinh ^{-1}(1)+36 \sinh ^{-1}\left(\frac{1}{\sqrt{2}}\right)\right)
$}
\end{equation*}
with $i_{B,12} \approx 0.04970171296728090$.

A non-trivial consistency check is supplied by the trace identity
\begin{equation*}
  3\,i_{B,11} \;=\; \int_{[0,1]^3}\!(1-u_1)(1-u_2)(1-u_3)\,\frac{u_1^2+u_2^2+u_3^2}{|\bm{u}|^3}\,d\bm{u} \;=\; i_A,
\end{equation*}
which is verified analytically and numerically to working precision.


\subsection{The 9{\times}9 body bilinear matrix on $T_9$}
\label{sec:results-body9}

With the master integrals of Sections~\ref{sec:results} and the three scalar building blocks $i_A$, $i_{B,11}$, $i_{B,12}$, the body bilinear form $\mathcal{B}_9 : T_9\times T_9 \to \mathbb{R}$ is assembled entry-by-entry from
\begin{equation}
  \mathcal{B}_9[i,j] \;=\; 32\,\sum_{p,q=1}^{3}\Bigl\{A\,\delta_{pq}\,K_{1}[\Pi_{ij,pq}] \,+\, B\,K_{3}[\Pi_{ij,pq};\,p,q]\Bigr\},
\end{equation}
where $\Pi_{ij,pq}(\bm{u})$ is the residual polynomial obtained by (a) substituting $\phi^{(i)}_p(\bm{r})\phi^{(j)}_q(\bm{r}')$, (b) integrating the centre-of-mass variable $\bm{\xi}$ over the overlap box, and (c) dividing out the standard tent factor $(1-u_1)(1-u_2)(1-u_3)$.  The prefactor $32 = 4\cdot 8/1$ absorbs the $s\mapsto 2u$ Jacobian $8$, the Green-tensor rescaling $G(2\bm{u})=\tfrac12 [A/|\bm{u}| + B\,u_p u_q/|\bm{u}|^3]$, and the symmetric-box factor $2/1$.  A further factor of $8$ appears on parity-even entries when we restrict to the positive octant $[0,1]^3$, so a diagonal displacement-block entry is proportional to $256$.

Under cubic symmetry the $T_9$-basis decomposes into four irreducible blocks, each of which is characterised by at most two independent scalars:
\begin{itemize}
  \item \textbf{Displacement block} ($i,j\in\{1,2,3\}$): diagonal, with a single scalar.
  \item \textbf{Displacement--strain block} ($i\in\{1,2,3\}$, $j\in\{4,\ldots,9\}$): identically zero by parity.
  \item \textbf{Axial-strain block} ($i,j\in\{4,5,6\}$): diagonal entry $\alpha$, off-diagonal entry $\beta$.
  \item \textbf{Axial--shear block} ($i\in\{4,5,6\}$, $j\in\{7,8,9\}$): identically zero by parity.
  \item \textbf{Shear-strain block} ($i,j\in\{7,8,9\}$): diagonal entry $\gamma$, off-diagonals identically zero by parity.
\end{itemize}

\subsubsection*{Displacement block}
All three displacement-block diagonal entries are equal and reduce in closed form to
\begin{equation}
\resizebox{\linewidth}{!}{$\displaystyle
  \mathcal{B}_9[1,1] \;=\; \mathcal{B}_9[2,2] \;=\; \mathcal{B}_9[3,3]
     \;=\; 256\,\bigl(A\,i_A \;+\; B\,i_{B,11}\bigr)
     \;=\; 60.23400462046913\,A \;+\; 20.07800154015638\,B,
$}
  \label{eq:B9disp}
\end{equation}

with off-diagonal entries vanishing by parity.

\subsubsection*{Axial-strain block}
The diagonal entry $\alpha := \mathcal{B}_9[4,4] = \mathcal{B}_9[5,5] = \mathcal{B}_9[6,6]$ splits into
\begin{equation}
\resizebox{\linewidth}{!}{$\displaystyle
  \alpha \;=\; \tfrac{256}{3}\Bigl[A\bigl(i_A - 2 K_1^{(2,0,0)}\bigr)
     \;+\; B\bigl(i_{B,11} - 2 K_3^{(2,0,0)}\bigr)\Bigr]
     \;=\; 15.37252824403423\,A \;+\; 3.949761457945144\,B,
$}
  \label{eq:alpha-axial}
\end{equation}
where
\begin{equation*}
\resizebox{\linewidth}{!}{$\displaystyle
  K_1^{(2,0,0)} \;:=\; \int_{[0,1]^3}(1-u_1)(1-u_2)(1-u_3)\,\frac{u_1^2}{|\bm{u}|}\,d\bm{u},
  \quad K_3^{(2,0,0)} \;:=\; \int_{[0,1]^3}(1-u_1)(1-u_2)(1-u_3)\,\frac{u_1^4}{|\bm{u}|^3}\,d\bm{u}.
$}
\end{equation*}
The off-diagonal entry $\beta := \mathcal{B}_9[4,5]$ is purely a $B$-channel correction,
\begin{equation}
  \beta \;=\; -256\,B\,K_{3,\text{off}}^{(1,1,0)}
     \;=\; 0\,A \;-\; 1.471925680511371\,B,
  \label{eq:beta-axial}
\end{equation}
where $K_{3,\text{off}}^{(1,1,0)} = \int_{[0,1]^3}(1-u_1)(1-u_2)(1-u_3)\,u_1 u_2\,\tfrac{u_1 u_2}{|\bm{u}|^3}\,d\bm{u}$.  Note that $\beta$ vanishes in the $A$-channel because the integrand of the $A$ piece is odd in $(u_1,u_2)$.

\subsubsection*{Shear-strain block}
The diagonal entry $\gamma := \mathcal{B}_9[7,7] = \mathcal{B}_9[8,8] = \mathcal{B}_9[9,9]$ evaluates to
\begin{equation}
\resizebox{\linewidth}{!}{$\displaystyle
  \gamma \;=\; \tfrac{128}{3}\Bigl[A\bigl(i_A - 2 K_1^{(2,0,0)}\bigr) + B\bigl(i_{B,11} - 2 K_3^{(0,2,0)}\bigr)\Bigr]
     \;-\; 128\,B\,K_{3,\text{off}}^{(1,1,0)}
     \;=\; 7.686264122017116\,A \;+\; 2.119728856266587\,B,
$}
  \label{eq:gamma-shear}
\end{equation}
where $K_3^{(0,2,0)} = \int_{[0,1]^3}(1-u_1)(1-u_2)(1-u_3)\,\tfrac{u_2^2 u_1^2}{|\bm{u}|^3}\,d\bm{u}$.  The relation $\gamma^{(A)} = \tfrac12\,\alpha^{(A)}$ is an exact identity: the $A$-channel of $\gamma$ is exactly half the $A$-channel of $\alpha$, coming from the bilinearity of the two-component shear basis $\phi^{(7)} = (0,r_3/2,r_2/2)$.

Numerical verification of cubic symmetry (all six independent identities) was performed to $10^{-20}$ precision; see the Mathematica script output.


\subsection{The 18{\times}18 quad--quad block on \texorpdfstring{$T_{27}$}{T27}}
\label{sec:results-body27quad}

Extending $T_9$ to $T_{27}$ adjoins the 18 quadratic vector fields
\begin{equation*}
  \phi^{(9 + 6(k-1) + m)}(\br) \;=\; q_m(\br)\,\be_k,
  \qquad k \in \{1,2,3\},\;\; m \in \{1,\ldots,6\},
\end{equation*}
where the six quadratic monomials in canonical Voigt order are
\begin{equation*}
  q_1 = r_1^2,\quad q_2 = r_2^2,\quad q_3 = r_3^2,\quad
  q_4 = r_2 r_3,\quad q_5 = r_1 r_3,\quad q_6 = r_1 r_2.
\end{equation*}
We refer to $q_{1..3}$ as the squared ($S$) modes and to $q_{4..6}$ as the cross ($X$) modes.  On the $18\times 18$ quad--quad sub-block, the body bilinear form decomposes into an $A$-channel and a $B$-channel,
\begin{equation}
  \mathcal{B}_{27}^{\text{quad}}[i,j]
  \;=\; 32\,\bigl(a_{ij}\,A \;+\; b_{ij}\,B\bigr),
  \label{eq:B27quad-split}
\end{equation}
where the scalars $a_{ij}, b_{ij}$ depend only on the orbit of $(i,j)$ under the combined O$_h$ action on the three coordinate axes and the $i\leftrightarrow j$ symmetry.

Enumerating all $18^2 = 324$ index pairs and partitioning by this action yields exactly \textbf{34 orbits}.  Of these, $\textbf{20 orbits}$ (total multiplicity $222$) are identically zero in both channels: each such vanishing is a parity cancellation forced by cubic symmetry (the integrand is odd in at least one coordinate axis under the $[-1,1]^3$ integration).  The remaining \textbf{14 orbits} (total multiplicity $102$) carry nontrivial values, and are listed in Table~\ref{tab:tier6-orbits}.

\begin{table}[h]
\centering\small
\resizebox{\linewidth}{!}{%
\begin{tabular}{llrrr}
\toprule
Representative $(i,j)$ & Structure & mult & $A$-value & $B$-value \\
\midrule
$(10,10)$ & $S$-$S$ diag, axis-matched   & $3$  & $17.126396135464$ & $ 7.086377754857$ \\
$(11,11)$ & $S$-$S$ diag, axis-unmatched & $6$  & $17.126396135464$ & $ 5.020009190303$ \\
$(10,11)$ & $S$-$S$ same-block off-diag, one matched & $12$ & $14.698679523965$ & $ 5.444954908314$ \\
$(11,12)$ & $S$-$S$ same-block off-diag, both unmatched & $6$  & $14.698679523965$ & $ 3.808769707336$ \\
$(13,13)$ & $X$-$X$ diag, no-factor-matched & $3$  & $ 3.988696969908$ & $ 1.644075434840$ \\
$(14,14)$ & $X$-$X$ diag, one-factor-matched & $6$  & $ 3.988696969908$ & $ 1.172310767534$ \\
\midrule
$(10,17)$ & $S$-$S$ cross-block (type Ia) & $6$  & $0$               & $ 0.654189191338$ \\
$(11,16)$ & $S$-$S$ cross-block (type Ib) & $6$  & $0$               & $ 0.654189191338$ \\
$(14,26)$ & $X$-$X$ cross-block (type I)  & $6$  & $0$               & $ 0.654189191338$ \\
$(13,20)$ & $X$-$X$ cross-block (type II) & $6$  & $0$               & $-0.389578544263$ \\
$(14,19)$ & $X$-$X$ cross-block (type III)& $6$  & $0$               & $-0.389578544263$ \\
$(10,21)$ & $S$-$X$ cross-block (type II) & $12$ & $0$               & $-0.216414577671$ \\
$(11,21)$ & $S$-$X$ cross-block (type III)& $12$ & $0$               & $-0.216414577671$ \\
\bottomrule
\end{tabular}%
}
\caption{The 13 non-vanishing orbits of the $18{\times}18$ quad--quad block $\mathcal{B}_{27}^{\text{quad}}$.  Indices $i,j$ are global ($10\le i,j\le 27$).  The $A$- and $B$-values in each row are the two numerical coefficients of $A$ and $B$ in \eqref{eq:B27quad-split}.  Multiplicities sum to $90$; the remaining $234$ of the $324$ entries sit in $21$ zero orbits forced by cubic parity.  Values are exact to all $12$ digits displayed; the underlying high-precision data (\texttt{WorkingPrecision}$\to 30$, see Section~\ref{sec:results-high-precision}) carries $\sim 28$ correct digits.}
\label{tab:tier6-orbits}
\end{table}

\paragraph{Independent scalars.}
The 13 non-zero orbits collapse to \textbf{3 distinct $A$-scalars and 9 distinct $B$-scalars}:
\begin{align*}
  \alpha_1 &\approx 17.126396135464, &
  \alpha_2 &\approx 14.698679523965, &
  \alpha_3 &\approx  3.988696969908, \\[4pt]
  \beta_1 &\approx  7.086377754857, &
  \beta_2 &\approx  5.020009190303, &
  \beta_3 &\approx  5.444954908314, \\
  \beta_4 &\approx  3.808769707336, &
  \beta_5 &\approx  1.644075434840, &
  \beta_6 &\approx  1.172310767534, \\
  \beta_8 &\approx  0.654189191338, &
  \beta_9 &\approx -0.389578544263, &
  \beta_{10} &\approx -0.216414577671.
\end{align*}
The numbering omits $\beta_7$ deliberately: the orbit $(11,17)$ that previously appeared to give $\beta_7 \approx 0.692769$ is, at \texttt{WorkingPrecision}$\to 30$, identically zero in both channels (see paragraph below and Section~\ref{sec:results-high-precision}).  Across all $324$ entries the observed within-orbit spread is $0$ exactly in the $A$-channel and $< 10^{-27}$ in the $B$-channel, so these 12 scalars completely determine the $18\times 18$ block.  All values are exported to \texttt{Mathematica/CubeT6ScalarValues\_HighPrec.wl}.

\paragraph{Symmetry structure and accidental identities.}
The $A$-channel vanishes on every cross-direction orbit ($k_i\neq k_j$) and on every $S$-$X$ orbit: this reflects the fact that the $A$ part of the elastic Green's tensor is proportional to $\delta_{ij}/|\mathbf{r}-\mathbf{r}'|$ (an isotropic trace piece), so it only connects modes carrying identical Cartesian vector components and with an even polynomial overlap.  The $B$-channel, which carries the $(\mathbf{r}{-}\mathbf{r}')_i(\mathbf{r}{-}\mathbf{r}')_j/|\mathbf{r}{-}\mathbf{r}'|^3$ anisotropic piece, is nonzero on all $14$ orbits.

An additional structural observation from Table~\ref{tab:tier6-orbits} is that the $7$ surviving cross-direction orbits collapse to only \emph{three} distinct $B$-values: $\beta_8=0.6542$ (appearing in three orbits), $\beta_9=-0.3896$ (appearing in two orbits), and $\beta_{10}=-0.2164$ (appearing in two orbits).  These collapses are identities to $\ge 27$-digit precision (\texttt{WorkingPrecision}$\to 30$) that are not a direct consequence of the O$_h$-orbit decomposition alone: our canonical labelling distinguishes $(10,17)$, $(11,16)$, and $(14,26)$ as separate O$_h$ orbits, yet they integrate to the same number.  Together with the unconditional vanishing of orbit $(11,17)$ (described next), they reflect a deeper structural identity of the cube Green-tensor double integral; closed-form verification via PSLQ on the high-precision values of Section~\ref{sec:results-high-precision} is the natural next step.

\paragraph{The vanishing of orbit $(11,17)$.}
A separate, qualitative finding of the high-precision recomputation is that the $S$-$X$ cross-block orbit at $(11,17)$ ($S$-$X$ type I, multiplicity $12$) is \emph{exactly zero in both channels}.  Earlier \texttt{MachinePrecision} assemblies returned $B \approx 0.692769$ for this orbit, which we previously labelled $\beta_7$; the overnight \texttt{WorkingPrecision}$\to 30$ recomputation collapses this value to $|B| < 10^{-27}$, consistent with an exact algebraic cancellation rather than residual noise.  We have therefore reclassified $(11,17)$ as a zero orbit and dropped $\beta_7$ from the list of independent scalars.  The previously reported "$14$ non-zero orbits, $13$ scalars" is replaced by the present "$13$ non-zero orbits, $12$ scalars".  This adjustment changes the multiplicity of zero orbits from $222$ to $234$ (i.e.\ $20+1$ orbits, mult $222+12$), and the multiplicity of non-zero orbits from $102$ to $90$.

\paragraph{Extraction method and precision.}
{\sloppy
The scalars $\alpha_i, \beta_j$ were initially computed by \texttt{Mathematica/\allowbreak CubeT6Block.wl}, which assembles $\mathcal{B}_{27}^{\text{quad}}$ via \texttt{NIntegrate} at \texttt{MachinePrecision} with \texttt{PrecisionGoal}\allowbreak$\to 12$.  That configuration completes all $324$ entries in $\sim 150$ s and cross-checks against the $T_9$ Tier-5 displacement diagonal \eqref{eq:B9disp} to better than $3\times 10^{-12}$, but it is just at the edge of double precision: the trailing $1$--$3$ decimal digits of the $\alpha_i,\beta_j$ are noise, and orbit $(11,17)$ in particular is misreported as having $B\approx 0.6928$ rather than its true value $0$.\par

The dedicated overnight high-precision recomputation \texttt{Mathematica/\allowbreak CubeT6OvernightHighPrec.wl} ran the same $14$ orbit representatives at \texttt{WorkingPrecision}\allowbreak$\to 30$ on $2026$-$04$-$08$, completing in $\sim 1217$ s of total \texttt{NIntegrate} time.  Its results are exported to \texttt{Mathematica/\allowbreak CubeT6ScalarValues\_\allowbreak HighPrec.wl} and tabulated in Section~\ref{sec:results-high-precision}.  The high-precision pass:\par}
\begin{enumerate}[label=(P\arabic*),leftmargin=2em]
\item confirms the $12$-digit values of the surviving $13$ orbits (corrections at the $11$th--$12$th decimal where machine-precision drift was visible);
\item shrinks the within-orbit spread of the $\alpha$- and $\beta$-channels to $< 10^{-27}$, raising the cubic-symmetry consistency check from $\sim 10^{-12}$ to $\sim 10^{-27}$;
\item promotes orbit $(11,17)$ from "$\beta_7 \approx 0.6928$" to a confirmed structural zero;
\item supplies the $\ge 28$ correct digits needed for PSLQ-based closed-form reconstruction of the $\alpha_i$ and $\beta_j$ in the same spirit as \eqref{eq:alpha-axial}--\eqref{eq:gamma-shear}.
\end{enumerate}
The 7 exact high-degree masters $\mathcal{M}^+_{(2,2,2)},\ldots,\mathcal{M}^+_{(8,0,0)}$ primed from \texttt{CubeT6Masters.wl} continue to anchor the leading parity-even contributions in closed form; PSLQ reconstruction of the remaining $\alpha_i,\beta_j$ in terms of $\sqrt{3}$, $\pi$, $\log$, $\tan^{-1}$, and $\sinh^{-1}$ atoms is the next analytical stage.


\subsection{High-precision scalar values}
\label{sec:results-high-precision}

{\sloppy
The values below are produced by \texttt{Mathematica/\allowbreak CubeT6OvernightHighPrec.wl} at \texttt{WorkingPrecision}\allowbreak$\to 30$ and stored verbatim in \texttt{Mathematica/\allowbreak CubeT6ScalarValues\_\allowbreak HighPrec.wl}.  Each entry carries $\ge 28$ correct decimal digits and is suitable as input to PSLQ-style integer-relation searches against an atom basis $\{1,\sqrt{3},\pi,\log 2,\log(1+\sqrt{3}),\log(\sqrt{3}-1),\sinh^{-1}1,\tan^{-1}(\cdot),\ldots\}$ in the spirit of Tables~\ref{tab:tier1-masters}--\ref{tab:tier2-masters}.\par}

\begin{table}[h]
\centering\small
\begin{tabular}{lll}
\toprule
Symbol & Value (28 digits) & Source orbit \\
\midrule
$\alpha_1$ & $\phantom{-}17.1263961354644384869032245660$ & $(10,10),(11,11)$ \\
$\alpha_2$ & $\phantom{-}14.6986795239648355731322989855$ & $(10,11),(11,12)$ \\
$\alpha_3$ & $\phantom{-1}3.9886969699081621215723702580$ & $(13,13),(14,14)$ \\
\midrule
$\beta_1$ & $\phantom{-1}7.0863777548574787932169236859$ & $(10,10)$ \\
$\beta_2$ & $\phantom{-1}5.0200091903034773113689403067$ & $(11,11)$ \\
$\beta_3$ & $\phantom{-1}5.4449549083142088200682904915$ & $(10,11)$ \\
$\beta_4$ & $\phantom{-1}3.8087697073364094814150175581$ & $(11,12)$ \\
$\beta_5$ & $\phantom{-1}1.6440754348395546266255471889$ & $(13,13)$ \\
$\beta_6$ & $\phantom{-1}1.1723107675343058603685866457$ & $(14,14)$ \\
$\beta_8$ & $\phantom{-1}0.6541891913383859981036052975$ & $(10,17),(11,16),(14,26)$ \\
$\beta_9$ & $-0.3895785442629410683383767348$ & $(13,20),(14,19)$ \\
$\beta_{10}$ & $-0.2164145776711156168035387731$ & $(10,21),(11,21)$ \\
\bottomrule
\end{tabular}
\caption{High-precision values of the $12$ independent scalars determining the $18\times 18$ quad--quad block of $\mathcal{B}_{27}^{\text{quad}}$, computed at \texttt{WorkingPrecision}$\to 30$.  The previously reported $\beta_7\approx 0.6928$ is absent: orbit $(11,17)$ has been re-classified as identically zero ($|B|<10^{-27}$).  Source orbits list every distinct $(i,j)$ representative whose entry equals the given scalar.}
\label{tab:high-prec-scalars}
\end{table}

\input{cube_galerkin27_closedforms.tex}

\paragraph{Per-orbit timing.}
The dominant cost of the overnight run was concentrated in two orbits: $(10,10)$ ($S$-$S$ axis-matched diagonal) at $989.6$ s, and $(10,11)$ ($S$-$S$ off-diagonal one-matched) at $140.3$ s.  All other orbits completed in $< 31$ s each, with the parity-zero $A$-channel orbits taking $< 15$ s and the smallest cross-block orbits in $\sim 10^{-3}$ s (cache hits).  Total wall time was $1217.07$ s, dominated by the single $(10,10)$ entry.

\paragraph{Verification.}
{\sloppy
Cubic-symmetry consistency for the 13 surviving orbits was checked at \texttt{WorkingPrecision}\allowbreak$\to 30$ and matches to $< 10^{-27}$ in all six independent identities.  The within-orbit spreads collapse from the $\sim 10^{-15}$ floating-point noise of the \texttt{MachinePrecision} run to $\sim 10^{-28}$, consistent with all displayed digits being algebraically exact.  Independent reproduction of $\alpha_1$ through $\alpha_3$ via the closed-form Tier-1 master integrals $\mathcal{M}_{p,q,r}$ (Table~\ref{tab:tier1-masters}) and the linear-combination assembly recipe of \eqref{eq:B27quad-split} agrees with the high-precision values to all $28$ digits, completing the cross-check.\par}


\subsection{Elastic stiffness bilinear form
\texorpdfstring{$\mathcal{B}_{\mathrm{el}}$}{B\_el} on
\texorpdfstring{$T_{27}$}{T\_27}}
\label{sec:results-Bel27}

Where the body bilinear form $\mathcal{B}_{\mathrm{body}}$ of Sections~\ref{sec:results-body9}--\ref{sec:results-body27quad} carries the (singular) Green-tensor double integral, the \emph{elastic} bilinear form $\mathcal{B}_{\mathrm{el}}$ is the ordinary strain-energy inner product of the isotropic stiffness contrast $\Delta c$ against the trial-space basis.  For the isotropic Kelvin form
\begin{equation}
  \Delta c_{ijkl}
  \;=\; \Delta\lambda\,\delta_{ij}\delta_{kl}
     \;+\; \Delta\mu\,\bigl(\delta_{ik}\delta_{jl} + \delta_{il}\delta_{jk}\bigr),
  \label{eq:deltaC-isotropic}
\end{equation}
the Galerkin matrix entries are the polynomial integrals
\begin{equation}
  \mathcal{B}_{\mathrm{el}}[v, w]
  \;=\; \int_{[-1,1]^3}\!\varepsilon(v):\Delta c:\varepsilon(w)\,dV
  \;=\; \int_{[-1,1]^3}\!\Bigl[\Delta\lambda\,\operatorname{tr}\varepsilon(v)\operatorname{tr}\varepsilon(w)
     \;+\; 2\Delta\mu\,\varepsilon(v){:}\varepsilon(w)\Bigr]\,dV,
  \label{eq:Bel-definition}
\end{equation}
with $\varepsilon(v)_{ij} = \tfrac12(\partial_i v_j + \partial_j v_i)$.  Because the $27$-dimensional trial space is polynomial, \eqref{eq:Bel-definition} is a closed-form symbolic integral and the full $27\times 27$ matrix $\mathcal{B}_{\mathrm{el}}$ is produced directly in \texttt{Mathematica/CubeT27Assemble.wl} Section~11 (no hand-coded block entries).

\paragraph{Structural result: the linear--quadratic cross block vanishes by parity.}
The decisive structural observation is that the $9\times 18$ block $\mathcal{B}_{\mathrm{el}}[\,1{:}9,\,10{:}27\,]$ connecting the linear subspace $T_9$ to the quadratic subspace $T_{27}^{\mathrm{quad}}$ is \emph{identically zero}:
\begin{equation}
  \mathcal{B}_{\mathrm{el}}\bigl[\,T_9,\,T_{27}^{\mathrm{quad}}\,\bigr]
  \;=\; 0.
  \label{eq:Bel-cross-zero}
\end{equation}
The mechanism is elementary.  The strain of any linear basis vector in $T_9$ is a \emph{constant} tensor, whereas the strain of any quadratic basis vector in $T_{27}^{\mathrm{quad}}$ is \emph{linear} in the coordinates $(r_1,r_2,r_3)$.  The integrand in \eqref{eq:Bel-definition} is therefore an odd polynomial on the symmetric box $[-1,1]^3$, and every component of every such cross entry integrates to $0$ term-by-term.  This cancellation is exact and algebraic, not numerical; the symbolic Mathematica run in \texttt{CubeT27Assemble.wl} confirms it by computing the $\sup$-norm of the full $9{\times}18$ block:
\begin{center}\texttt{|BelSym[[1..9,\,10..27]]|\ max\ =\ 0}.\end{center}

The immediate Galerkin-closure consequence is that $\mathcal{B}_{\mathrm{el}}$ is \emph{block-diagonal} with respect to the decomposition $T_{27} = T_9 \oplus T_{27}^{\mathrm{quad}}$:
\begin{equation}
  \mathcal{B}_{\mathrm{el}}
  \;=\;
  \begin{pmatrix}
    \mathcal{B}_{\mathrm{el}}^{\,T_9}   & 0 \\[2pt]
    0 & \mathcal{B}_{\mathrm{el}}^{\,\text{quad--quad}}
  \end{pmatrix}.
  \label{eq:Bel-block-diagonal}
\end{equation}
This stands in sharp contrast to the body form $\mathcal{B}_{\mathrm{body}}$, whose cross block is \emph{not} zero (see the Section~\ref{sec:results-body27quad} YY block and the $T_9{\times} T_{27}^{\mathrm{quad}}$ couplings that drive the body-channel Schur complement).

\paragraph{The $T_9$ local block.}
On the $9{\times}9$ linear sub-block, constants contribute no strain and the $3{\times}3$ displacement block is identically zero.  The remaining $6{\times}6$ strain sub-block (rows/cols $4\ldots 9$) reproduces the classical local Kelvin form $V\cdot\varepsilon{:}\Delta c{:}\varepsilon$ with $V=8$:
\begin{equation}
  \mathcal{B}_{\mathrm{el}}^{\,T_9}\bigl[\,4{:}9,\,4{:}9\,\bigr]
  \;=\; 8
  \begin{pmatrix}
    \Delta\lambda+2\Delta\mu & \Delta\lambda & \Delta\lambda & 0 & 0 & 0\\
    \Delta\lambda & \Delta\lambda+2\Delta\mu & \Delta\lambda & 0 & 0 & 0\\
    \Delta\lambda & \Delta\lambda & \Delta\lambda+2\Delta\mu & 0 & 0 & 0\\
    0 & 0 & 0 & \Delta\mu & 0 & 0\\
    0 & 0 & 0 & 0 & \Delta\mu & 0\\
    0 & 0 & 0 & 0 & 0 & \Delta\mu
  \end{pmatrix}.
  \label{eq:Bel-T9-strain}
\end{equation}
The $\mathrm{O}_h$-irreducible eigenvalues of \eqref{eq:Bel-T9-strain} are
\begin{align}
  \mathcal{B}_{\mathrm{el}}^{\,T_9}\Big|_{A_{1g}} \;&=\; 8\bigl(3\Delta\lambda + 2\Delta\mu\bigr)
  \;=\; 8\cdot 3K_{\Delta}, \label{eq:Bel-T9-A1g}\\
  \mathcal{B}_{\mathrm{el}}^{\,T_9}\Big|_{E_{g}}   \;&=\; 16\,\Delta\mu,                       \label{eq:Bel-T9-Eg}\\
  \mathcal{B}_{\mathrm{el}}^{\,T_9}\Big|_{T_{2g}}  \;&=\; \phantom{1}8\,\Delta\mu,             \label{eq:Bel-T9-T2g}
\end{align}
where $K_{\Delta} := \Delta\lambda + \tfrac23\Delta\mu$ is the bulk contrast.  The $T_9$ strain block carries exactly \emph{two} independent scalars $(\Delta\lambda,\Delta\mu)$ -- it is fully isotropic -- whereas the Path-A coefficients $(T_{1c},T_{2c},T_{3c})$ of the cube effective $T$-matrix carry \emph{three} independent scalars, the third one being the cubic-anisotropy piece associated with the static Eshelby integral $C$ (see the Path-A reference in \texttt{cubic\_scattering/\allowbreak effective\_contrasts.py}).  Equations \eqref{eq:Bel-T9-A1g}--\eqref{eq:Bel-T9-T2g} make it immediate that the missing cubic anisotropy cannot come from the stiffness channel in isolation.

\paragraph{The quad--quad $18{\times}18$ block.}
On the $18{\times}18$ sub-block $\mathcal{B}_{\mathrm{el}}^{\,\text{quad--quad}}$ with the \texttt{basisT27quad} ordering of Section~\ref{sec:results-body27quad}, the symbolic integration yields a matrix with $54$ non-zero entries out of $324$ (structurally sparse, diagonal-dominant) and symbolic rank $18$.  Per-direction the $6{\times}6$ diagonal splits into four distinct scalar entries, reflecting whether the quadratic monomial is of $S$ or $X$ type and whether its factor axes match the direction index $k$:
\begin{equation}
  \operatorname{diag}\!\bigl(\mathcal{B}_{\mathrm{el}}^{\,\text{quad--quad}}\bigr)_{\text{dir.}\,k}
  \;=\;
  \begin{cases}
    \tfrac{32}{3}\bigl(\Delta\lambda + 2\Delta\mu\bigr)
      & q = r_k^2 \quad\text{($S$, matched)},\\[6pt]
    \tfrac{32}{3}\,\Delta\mu
      & q = r_m^2,\ m\neq k \quad\text{($S$, unmatched)},\\[6pt]
    \tfrac{16}{3}\,\Delta\mu
      & q = r_m r_n,\ m,n\neq k \quad\text{($X$, no match)},\\[6pt]
    \tfrac{8}{3}\bigl(\Delta\lambda + 3\Delta\mu\bigr)
      & q = r_k r_m,\ m\neq k \quad\text{($X$, one match)}.
  \end{cases}
  \label{eq:Bel-quad-diagonal}
\end{equation}
The four per-direction diagonals \eqref{eq:Bel-quad-diagonal} are the analogues, on the quadratic sub-basis, of the three $T_9$ eigenvalues \eqref{eq:Bel-T9-A1g}--\eqref{eq:Bel-T9-T2g}.  Their orbit structure under $\mathrm{O}_h$ matches exactly the parity classification of Table~\ref{tab:tier6-orbits}.

\paragraph{Consequence for the Schur-complement reduction to $T_9$.}
Under the body-channel Galerkin eigenproblem $(\mathcal{M}_{27} - \varepsilon\,\mathcal{B})\,\mathbf{c} = 0$ the Schur-complemented effective operator on $T_9$ is
\begin{equation*}
  T_9^{\mathrm{eff}}
  \;=\; \bigl(\mathcal{M}_{11} - \mathcal{M}_{12}\mathcal{M}_{22}^{-1}\mathcal{M}_{21}\bigr)^{-1}
     \bigl(\mathcal{B}_{11} \;-\; \mathcal{M}_{12}\mathcal{M}_{22}^{-1}\mathcal{B}_{21} \;-\; \mathcal{B}_{12}\mathcal{M}_{22}^{-1}\mathcal{M}_{21}\bigr).
\end{equation*}
For the stiffness channel, $\mathcal{B} = \mathcal{B}_{\mathrm{el}}$ and \eqref{eq:Bel-cross-zero} forces $\mathcal{B}_{12} = \mathcal{B}_{21}^{\top} = 0$, so the middle and trailing terms vanish identically and the Schur-complemented stiffness operator reduces to the local block \eqref{eq:Bel-T9-strain}:
\begin{equation}
  T_9^{\mathrm{eff,\,el}}
  \;=\; \bigl(\mathcal{M}_{11} - \mathcal{M}_{12}\mathcal{M}_{22}^{-1}\mathcal{M}_{21}\bigr)^{-1}\,
     \mathcal{B}_{\mathrm{el}}^{\,T_9}\bigl[\,1{:}9,\,1{:}9\,\bigr].
  \label{eq:T9-eff-stiffness-schur}
\end{equation}
In particular, the entire $T_{27}^{\mathrm{quad}}$ sub-block of $\mathcal{B}_{\mathrm{el}}$ drops out of the $T_9$ reduction at leading order, and the cubic-anisotropy scalar $T_{3c}$ of the Path-A effective $T$-matrix \emph{cannot} originate in the elastic Galerkin channel on the quadratic sub-basis.  The parallel observation of Section~\ref{sec:results-body9} was that the body-channel $T_9^{\mathrm{eff}}$ acts as a pure scalar on the $6$-dim strain sector in the numerical $A$-polarization probe of Section~12 (Schur-complemented diagonals \texttt{const = -47.1737...}, \texttt{strain = 5.7647...}, off-diagonals exactly $0$ at machine precision).  That isolated probe left the origin of $T_{3c}$ ambiguous: it might have been a body-only effect visible only in the second polarization channel, or a genuinely mixed body--stiffness quantity.  The symbolic computation in Section~12b below resolves this ambiguity by keeping both polarization channels explicit, and shows that $T_{3c}$ arises from the $B$-polarization of the body operator alone (see \eqref{eq:T3c-signature}); the stiffness channel contributes only isotropic corrections.

\subsubsection*{Symbolic verification (Section 11b)}
\label{sec:results-bel-schur-11b}
An explicit symbolic computation of the stiffness Schur complement has been added as Section~11b of \texttt{CubeT27Assemble.wl} and run end-to-end at exact rational precision.  Partitioning the $27{\times}27$ matrices into $9\mid 18$ blocks
\[
  \mathcal{M}_{27} = \begin{pmatrix}\mathcal{M}_{11} & \mathcal{M}_{12}\\ \mathcal{M}_{21} & \mathcal{M}_{22}\end{pmatrix},
  \qquad
  \mathcal{B}_{\mathrm{el}} = \begin{pmatrix}\mathcal{B}_{11}^{\mathrm{el}} & \mathcal{B}_{12}^{\mathrm{el}}\\ \mathcal{B}_{21}^{\mathrm{el}} & \mathcal{B}_{22}^{\mathrm{el}}\end{pmatrix},
\]
and forming the full Schur complement
\[
  \mathcal{B}_{\mathrm{el}}^{\mathrm{Schur}}
  \;\equiv\; \mathcal{B}_{11}^{\mathrm{el}}
     - \mathcal{M}_{12}\mathcal{M}_{22}^{-1}\mathcal{B}_{21}^{\mathrm{el}}
     - \mathcal{B}_{12}^{\mathrm{el}}\mathcal{M}_{22}^{-1}\mathcal{M}_{21},
\]
Mathematica returns $\mathsf{Simplify}\bigl[\mathcal{B}_{\mathrm{el}}^{\mathrm{Schur}} - \mathcal{B}_{11}^{\mathrm{el}}\bigr] = 0$ identically (not just numerically), confirming \emph{algebraically} that the parity-vanishing result $\mathcal{B}_{12}^{\mathrm{el}} = (\mathcal{B}_{21}^{\mathrm{el}})^{\top} = 0$ makes the middle and trailing Schur terms disappear exactly.  The strain-sector (rows/cols $4\dots 9$) eigenvalues of $\mathcal{B}_{\mathrm{el}}^{\mathrm{Schur}}$ are then the pure isotropic Kelvin form,
\begin{equation}
  \lambda_{A_{1g}} \;=\; 8\,(3\,\Delta\lambda + 2\,\Delta\mu),
  \qquad
  \lambda_{E_{g}} \;=\; 16\,\Delta\mu,
  \qquad
  \lambda_{T_{2g}} \;=\; 8\,\Delta\mu,
  \label{eq:Bel-Schur-eigs}
\end{equation}
matching the direct diagonal extraction \eqref{eq:Bel-T9-A1g}--\eqref{eq:Bel-T9-T2g} of Section~11.  Post-multiplying by the body-channel mass Schur inverse to form the full stiffness-channel effective operator on $T_9$,
\[
  T_9^{\mathrm{eff,\,el}}
  \;=\; \bigl(\mathcal{M}_{11} - \mathcal{M}_{12}\mathcal{M}_{22}^{-1}\mathcal{M}_{21}\bigr)^{-1}\,\mathcal{B}_{\mathrm{el}}^{\mathrm{Schur}},
\]
the symbolic simplification gives strain-sector irrep eigenvalues
\begin{equation}
  \lambda_{A_{1g}}\bigl(T_9^{\mathrm{eff,\,el}}\bigr) = 9\,\Delta\lambda + 6\,\Delta\mu,
  \qquad
  \lambda_{E_{g}}\bigl(T_9^{\mathrm{eff,\,el}}\bigr) = 6\,\Delta\mu,
  \qquad
  \lambda_{T_{2g}}\bigl(T_9^{\mathrm{eff,\,el}}\bigr) = 6\,\Delta\mu.
  \label{eq:T9-eff-el-eigs}
\end{equation}
The crucial observation is the equality
\begin{equation}
  \lambda_{E_g}\bigl(T_9^{\mathrm{eff,\,el}}\bigr)
  \;=\; \lambda_{T_{2g}}\bigl(T_9^{\mathrm{eff,\,el}}\bigr)
  \;=\; 6\,\Delta\mu,
  \label{eq:Eg-T2g-degeneracy}
\end{equation}
which holds \emph{symbolically} in $\Delta\lambda, \Delta\mu$ with no residual.  Cubic anisotropy, by definition, requires $\lambda_{E_g} \neq \lambda_{T_{2g}}$ on the five-dimensional deviatoric strain sector (the $E_g \oplus T_{2g}$ splitting); the fact that they are \emph{exactly equal} in $T_9^{\mathrm{eff,\,el}}$ therefore rules out any cubic-anisotropic component from the stiffness channel in isolation.  Equivalently, $T_9^{\mathrm{eff,\,el}}$ is strictly proportional to the isotropic Kelvin contraction, and its Path-A decomposition acquires no contribution to $T_{3c}$.  The complementary question -- whether the body channel by itself already produces a nonzero $T_{3c}$, or whether cross-coupling to $\mathcal{B}_{\mathrm{el}}$ is required -- is answered symbolically in Section~12b below.

\subsubsection*{Origin of $T_{3c}$: combined body+stiffness Schur (Section 12b)}
\label{sec:results-body-stiffness-12b}

With the stiffness channel definitively isotropic \eqref{eq:Eg-T2g-degeneracy}, the origin of the cube's $T_{3c}$ cubic-anisotropy scalar must lie in the body channel.  Section~12b of \texttt{CubeT27Assemble.wl} completes this identification by symbolically constructing the combined body+stiffness Schur-complemented effective operator
\begin{equation}
  T_9^{\mathrm{eff,\,tot}}
  \;=\; \bigl(\mathcal{M}_{11}
     - \mathcal{M}_{12}\mathcal{M}_{22}^{-1}\mathcal{M}_{21}\bigr)^{-1}
     \Bigl[\eta_{\mathrm{body}}\,\mathcal{B}_{\mathrm{body}}^{\mathrm{Schur}}
       + \mathcal{B}_{\mathrm{el}}^{\mathrm{Schur}}\Bigr],
  \label{eq:T9-eff-total-schur}
\end{equation}
where $\eta_{\mathrm{body}}$ is a free body-coupling coefficient (physically $\propto \omega^{2}\Delta\rho$) and $\mathcal{B}_{\mathrm{body}}^{\mathrm{Schur}}$ is the $9{\times}9$ Schur complement of the full $27{\times}27$ body bilinear form.  The key technical step is to keep $\mathcal{B}_{\mathrm{body}}^{\mathrm{Schur}}$ symbolic in the two independent polarization channels $(A_{\mathrm{elas}}, B_{\mathrm{elas}})$ of the underlying elastic Green's tensor:
\begin{itemize}
  \item The $A$-channel (coefficient $A_{\mathrm{elas}}$) collects the isotropic $\delta_{ij}$ piece of the Green's tensor, weighted by the $K_{1}$ kernel (weak $1/|r|$ integrals over the Mp masters of Tier-6).
  \item The $B$-channel (coefficient $B_{\mathrm{elas}}$) collects the tensorial anisotropic piece, i.e.\ the $r_{i}r_{j}/|r|^{3}$ contraction weighted by the $K_{3}$ kernel (DuffyCoordinates octant integrals).
\end{itemize}
These correspond to the two independent building blocks $(\mathsf{aPiece},\mathsf{bPiece})$ returned by \texttt{bbBodyAB} in \texttt{CubeT6Block.wl}\@: the $A$-piece is non-zero only when the outer polarization indices $p=q$ (scalar projector), while the $B$-piece carries the full tensor structure through the $K_{3}$ kernel.

Substituting the Pell-simplified atomic scalars (\texttt{atomRules}) while holding $(A_{\mathrm{elas}}, B_{\mathrm{elas}})$ symbolic, Mathematica returns for $\mathcal{B}_{\mathrm{body}}^{\mathrm{Schur}}$ restricted to the strain sector (rows/cols $4\dots 9$) the irrep eigenvalues
\begin{equation}
\begin{aligned}
  \lambda_{A_{1g}}\bigl(\mathcal{B}_{\mathrm{body}}^{\mathrm{Schur}}\bigr) &= 15.3725\,A_{\mathrm{elas}} + 1.0059\,B_{\mathrm{elas}}, \\
  \lambda_{E_{g}}\bigl(\mathcal{B}_{\mathrm{body}}^{\mathrm{Schur}}\bigr) &= 15.3725\,A_{\mathrm{elas}} + 5.4217\,B_{\mathrm{elas}}, \\
  \lambda_{T_{2g}}\bigl(\mathcal{B}_{\mathrm{body}}^{\mathrm{Schur}}\bigr) &= \phantom{0}7.6863\,A_{\mathrm{elas}} + 2.1197\,B_{\mathrm{elas}},
\end{aligned}
\label{eq:Bbody-Schur-eigs}
\end{equation}
all to $\sim\!16$-digit MachinePrecision.  Both channels exhibit $\lambda_{E_g} \neq \lambda_{T_{2g}}$ at this raw Schur-complement stage, so a naive cubic-anisotropy test on $\mathcal{B}_{\mathrm{body}}^{\mathrm{Schur}}$ alone would (misleadingly) flag both channels.  However, the physically relevant operator is the full effective $T_9^{\mathrm{eff,\,tot}}$ obtained after pre-multiplying by the mass Schur inverse.  The symbolic simplification gives
\begin{equation}
\begin{aligned}
  \lambda_{A_{1g}}\bigl(T_9^{\mathrm{eff,\,tot}}\bigr) &= 9\,\Delta\lambda + 6\,\Delta\mu + \bigl(5.7647\,A_{\mathrm{elas}} + 0.3772\,B_{\mathrm{elas}}\bigr)\,\eta_{\mathrm{body}}, \\
  \lambda_{E_{g}}\bigl(T_9^{\mathrm{eff,\,tot}}\bigr) &= \phantom{9\,\Delta\lambda + {}}6\,\Delta\mu + \bigl(5.7647\,A_{\mathrm{elas}} + 2.0331\,B_{\mathrm{elas}}\bigr)\,\eta_{\mathrm{body}}, \\
  \lambda_{T_{2g}}\bigl(T_9^{\mathrm{eff,\,tot}}\bigr) &= \phantom{9\,\Delta\lambda + {}}6\,\Delta\mu + \bigl(5.7647\,A_{\mathrm{elas}} + 1.5898\,B_{\mathrm{elas}}\bigr)\,\eta_{\mathrm{body}}.
\end{aligned}
\label{eq:T9-eff-total-eigs}
\end{equation}
The $A$-channel coefficients on $\eta_{\mathrm{body}}$ are now \emph{all equal} to $5.7647$: the anisotropy visible in \eqref{eq:Bbody-Schur-eigs} in the $A$-polarization cancels exactly against the equivariant mass-Schur inverse, reproducing isotropy.  In the $B$-channel the coefficients split as $(0.3772,\,2.0331,\,1.5898)$, yielding the cubic-anisotropy signature
\begin{equation}
  \lambda_{E_{g}}\bigl(T_9^{\mathrm{eff,\,tot}}\bigr) - \lambda_{T_{2g}}\bigl(T_9^{\mathrm{eff,\,tot}}\bigr)
  \;=\; 0.4433\,B_{\mathrm{elas}}\,\eta_{\mathrm{body}}
  \;\neq\; 0,
  \label{eq:T3c-signature}
\end{equation}
which is \emph{independent} of both the stiffness contrasts $(\Delta\lambda, \Delta\mu)$ and the $A$-polarization of the body operator.

Equation \eqref{eq:T3c-signature} is the positive identification of the cube's $T_{3c}$ anisotropy at leading order in the $T_{27}$ Galerkin closure.  Three consequences deserve emphasis:
\begin{enumerate}
  \item The ``stiffness-only'' probe $\eta_{\mathrm{body}} = 0$ gives $\lambda_{E_g} - \lambda_{T_{2g}} = 0$ (Section~11b result), confirming no cubic-anisotropy component from $\mathcal{B}_{\mathrm{el}}$ alone.
  \item The ``$A$-only body'' probe $B_{\mathrm{elas}} = 0$ also gives $\lambda_{E_g} - \lambda_{T_{2g}} = 0$: the trace/isotropic piece $\delta_{ij}$ of the elastic Green's tensor (K$_{1}$ kernel) is insufficient, even when fully Galerkin-coupled through $T_{27}^{\mathrm{quad}}$.
  \item The $B$-channel $r_{i}r_{j}/|r|^{3}$ tensor piece of the elastic Green's tensor (K$_{3}$ kernel) is the sole source of the leading-order $T_{3c}$ signature in the tier-$27$ closure.
\end{enumerate}

This refines the preliminary expectation (preceding subsection) that $T_{3c}$ was an intrinsically mixed body--stiffness quantity: Section~12b establishes that the $B$-polarization of the body channel alone already breaks the $E_{g} / T_{2g}$ degeneracy, and the stiffness channel contributes only isotropic corrections.  The numerical coefficient $0.4433\,B_{\mathrm{elas}}\,\eta_{\mathrm{body}}$ is a closed form awaiting PSLQ/Pell simplification against the same basis of elementary constants used for the Tier-6 scalars $\{\alpha_{i}\}_{i=1}^{3}$ and $\{\beta_{j}\}_{j=1}^{9}$ of \texttt{CubeT6ScalarValues\_HighPrec\_Pell.wl}; this final symbolic simplification is the natural next step for a closed-form expression of the cube's $T_{3c}$ at leading Galerkin order.



%======================================================================
\section{Validation plan: Path-A $\Leftrightarrow$ Path-B for the 9-component subsystem}
%======================================================================

The existing \texttt{CubeAnalytic.wl} package implements Path A with
analytical Eshelby-delta corrections (the $(5\sqrt{3} - 2\pi)$ block from
the old \texttt{GreensTensorScatteringPaper\_Old.nb}).  As a correctness
test for the framework presented here, we will:
\begin{enumerate}[label=(V\arabic*)]
\item Compute the Path-B mass and stiffness matrices on the 9-dimensional
      subspace $T_9 \subset T_{27}$ spanned by constants and symmetric
      linear polynomials.
\item Solve the resulting $9\times 9$ generalized eigenproblem for the
      effective contrasts $\Delta\rho^{*}, \Delta\lambda^{*}, \Delta\mu^{*}$.
\item {\sloppy Compare to the Path-A values produced by
      \texttt{compute\_\allowbreak cube\_\allowbreak tmatrix} in
      \texttt{cubic\_\allowbreak scattering/\allowbreak effective\_\allowbreak contrasts.py}.\par}
\end{enumerate}
Path A and Path B do not need to give identical matrix entries (they
correspond to different finite-dimensional projections of the same
operator), but the effective contrasts must agree to leading order in
the polynomial truncation.  Any discrepancy quantifies the truncation
error and bounds the systematic error of Path A.

%======================================================================
\section{New integrals required for the 27-component closure}
%======================================================================

Beyond the constants used in the 9-component closure
($\mathcal{I}_{0,0,0}$, $\mathcal{J}^{(11)}_{0,0,0}$,
$\mathcal{J}^{(11)}_{2,0,0}$, etc.), the $27$-component closure
introduces master integrals with total monomial degree up to $5$.
The new families, organized by parity-allowed multi-indices and modulo
cubic symmetry, are:

\begin{table}[h]
\centering
\begin{tabular}{lll}
\toprule
$(p,q,r)$ class & rank & needed for \\
\midrule
$(2,2,0),\;(4,0,0)$         & 4 & $\langle \bphi^{(\text{lin})},G\bphi^{(\text{lin})}\rangle$ block \\
$(2,2,2),\;(4,2,0),\;(6,0,0)$ & 6 & $\langle \bphi^{(\text{quad})},G\bphi^{(\text{lin})}\rangle$ block \\
$(4,2,2),\;(4,4,0),\;(6,2,0),\;(8,0,0)$ & 8 & $\langle \bphi^{(\text{quad})},G\bphi^{(\text{quad})}\rangle$ block \\
\bottomrule
\end{tabular}
\caption{New cubic-symmetric master integral families needed for the
27-component Path-B closure. Each row corresponds to one block of the
$27\times 27$ stiffness matrix.}
\label{tab:new-masters}
\end{table}

These integrals will be derived analytically in Section~\ref{sec:tier1}
and the subsequent stages of \texttt{CubeGalerkin27.wl}.

%======================================================================
\section{Closed-form T-matrix via irrep decomposition}
\label{sec:closedform-tmatrix}
%======================================================================

The Lippmann--Schwinger Galerkin equation on $T_{27}$ contains two
channels:
\begin{itemize}
\item a \emph{body} bilinear form $B_{\text{body}}$
  (Sections~5--7, linear in the Kelvin coefficients $a_0, b_0$),
  from the $\omega^2\D\rho$ body-force term; and
\item an \emph{elastic-contrast} bilinear form from the $\D\mathbf{c}$
  stiffness perturbation, which enters through the \emph{Eshelby
  tensor}~$\mathbf{S}$---the volume integral of the second derivative
  of the Green's tensor $\int\partial^2 G_{ij}/\partial x_k\partial x_l\,dV$.
\end{itemize}

Under $O_h$ symmetry, the two parity sectors decouple: the
\emph{gerade} (German ``even''; subscript~$g$) irreps have basis
functions that are symmetric under spatial inversion
$\mathbf{x} \to -\mathbf{x}$, while the \emph{ungerade} (``odd'';
subscript~$u$) irreps are antisymmetric.  Physically, the gerade
irreps carry the strain modes and the ungerade irreps carry the
displacement and quadratic modes.  Each sector sees only one channel
at leading order:
\begin{itemize}
\item \textbf{Ungerade}: the body bilinear form
  $\varepsilon\,B_{\text{body}}$ with $\varepsilon = \omega^2\D\rho$.
\item \textbf{Gerade}: the LS-convolved stiffness
  $b^{\text{LS}}_\rho = \langle\boldsymbol{\varepsilon}_\rho
  \,|\,\mathbf{S}:\D\mathbf{c}\,|\,\boldsymbol{\varepsilon}_\rho\rangle$,
  obtained from the Eshelby integral $A^c, B^c, C^c$.
\end{itemize}

\paragraph{Dimensional consistency.}
The mass matrix $M_\rho$ and the LS-convolved bilinear forms
$\varepsilon b_\rho$ (body) and $b^{\text{LS}}_\rho$ (stiffness) are
all dimensionless and $O(1)$ in natural units, because the Green's
tensor carries the factor $1/(\rho v^2)$ that non-dimensionalizes
the contrast.  The \emph{direct} strain-energy integral
$\beta_\rho = \int\boldsymbol{\varepsilon}:\D\mathbf{c}:\boldsymbol{\varepsilon}\,dV$
has dimensions of~Pa and is $O(\D\mu)$, making it incompatible with~$M_\rho$
in the per-irrep formula.  Only the LS-convolved form may be combined
with the mass matrix.

\subsection{Cubic symmetry decomposition}

Under the cubic group $O_h$, the representation on $T_{27}$ decomposes as
\begin{equation}
  T_{27} = 4\,T_{1u} \oplus 2\,T_{2u} \oplus A_{2u} \oplus E_u
           \oplus A_{1g} \oplus E_g \oplus T_{2g}.
  \label{eq:Oh-decomp}
\end{equation}
All bilinear forms are $O_h$-equivariant, so Schur's lemma guarantees
block-diagonalization in the symmetry-adapted basis $U_\rho$
(Section~9c).  Each irrep $\rho$ with multiplicity $m_\rho$ contributes
an independent $m_\rho \times m_\rho$ block.

The ungerade and gerade sectors have different leading-order structure:
\begin{itemize}
\item \textbf{Ungerade} ($T_{1u}$, $T_{2u}$, $A_{2u}$, $E_u$):
  T-matrix dominated by the body bilinear form.
  $T_\rho = \varepsilon\,B_{\text{body},\rho}\,
  (M_\rho - \varepsilon\,B_{\text{body},\rho})^{-1}$.
\item \textbf{Gerade} ($A_{1g}$, $E_g$, $T_{2g}$):
  T-matrix determined by the LS-convolved stiffness with
  self-consistent amplification (Section~\ref{sec:strain-sector}).
\end{itemize}

\begin{table}[h]
\centering
\begin{tabular}{lcccl}
\toprule
Irrep & $d_\rho$ & $m_\rho$ & Block size & Physical content \\
\midrule
$T_{1u}$ & 3 & 4 & $4\times 4$ & displacement $+$ 3 quadratic modes \\
$T_{2u}$ & 3 & 2 & $2\times 2$ & 2 quadratic modes \\
$A_{2u}$ & 1 & 1 & $1\times 1$ & quadratic (pure) \\
$E_u$    & 2 & 1 & $1\times 1$ & quadratic (pure) \\
$A_{1g}$ & 1 & 1 & $1\times 1$ & volumetric strain \\
$E_g$    & 2 & 1 & $1\times 1$ & deviatoric axial strain \\
$T_{2g}$ & 3 & 1 & $1\times 1$ & off-diagonal shear strain \\
\bottomrule
\end{tabular}
\caption{The seven irrep blocks of the $27\times 27$ T-matrix.
The total dimension is $\sum d_\rho m_\rho = 12 + 6 + 1 + 2 + 1 + 2 + 3 = 27$.}
\label{tab:irrep-blocks}
\end{table}

\subsection{Strain-sector scalars (gerade irreps)}
\label{sec:strain-sector}

The three gerade irreps ($A_{1g}$, $E_g$, $T_{2g}$) are all
$m_\rho = 1$, yielding scalar T-matrix blocks.  Only the
LS-convolved stiffness contributes (the body force couples to
ungerade modes, not strain modes).  Self-consistent amplification
(Gubernatis et al.\ 1977) accounts for the scatterer modifying its
own internal field, coupling the bulk ($A_{1g}$) and deviatoric
($E_g$) channels through the effective~$\D\lambda^*$ formula.

\paragraph{Step 1: Born-level Eshelby contraction.}
The Eshelby integrals $A^c$, $B^c$, $C^c$ contracted with the
\emph{bare} stiffness contrast $\D\mathbf{c}$ yield Born-level
scalars $T_1^{c,0}$, $T_2^{c,0}$, $T_3^{c,0}$ and Born-level
per-irrep eigenvalues:
\begin{equation}
  \sigma_{A_{1g}}^0 = 3T_1^{c,0} + 2T_2^{c,0} + T_3^{c,0}, \quad
  \sigma_{E_g}^0 = 2T_2^{c,0} + T_3^{c,0}, \quad
  \sigma_{T_{2g}}^0 = 2T_2^{c,0}.
  \label{eq:born-sigma-irrep}
\end{equation}

\paragraph{Step 2: Amplification factors.}
Each irrep yields a self-consistent amplification factor:
\begin{align}
  A_\theta     &= (1 - \sigma_{A_{1g}}^0)^{-1}  &&\text{(volumetric)}, \\
  A_e^{\text{diag}} &= (1 - \sigma_{E_g}^0)^{-1} &&\text{(deviatoric diagonal)}, \\
  A_e^{\text{off}}  &= (1 - \sigma_{T_{2g}}^0)^{-1} &&\text{(off-diagonal shear)}.
\end{align}

\paragraph{Step 3: Effective contrasts.}
The amplified stiffness has \emph{cubic symmetry} even when the bare
contrast is isotropic, because $A_e^{\text{off}} \neq A_e^{\text{diag}}$
whenever $T_3^c \neq 0$ (i.e.\ whenever the scatterer is not a sphere):
\begin{align}
  \D\lambda^* &= \big(\D\lambda + \tfrac{2}{3}\D\mu\big)\,A_\theta
    - \tfrac{2}{3}\D\mu\,A_e^{\text{diag}},
    \label{eq:Dlambda-star} \\
  \D\mu^*_{\text{off}}  &= \D\mu \cdot A_e^{\text{off}}, \\
  \D\mu^*_{\text{diag}} &= \D\mu \cdot A_e^{\text{diag}}.
\end{align}
The cross-term in~\eqref{eq:Dlambda-star} couples the volumetric
and deviatoric amplification factors---this is the $\kappa$--$\mu$
coupling that a na\"ive per-irrep resummation would miss.

\paragraph{Step 4: Amplified T-matrix.}
The effective stiffness $\D\mathbf{c}^*$ decomposes into an
isotropic part $(\D\lambda^*, \D\mu^*_{\text{off}})$ plus a cubic
correction $2(\D\mu^*_{\text{diag}} - \D\mu^*_{\text{off}})\,E_{ijkl}$:
\begin{align}
  T_1^c &= T_1^{c,\text{iso}} + 2\delta\mu^*\, B^c, \label{eq:T1c-amp} \\
  T_2^c &= T_2^{c,\text{iso}}, \label{eq:T2c-amp} \\
  T_3^c &= T_3^{c,\text{iso}} + 2\delta\mu^*\,(A^c + B^c + C^c),
  \label{eq:T3c-amp}
\end{align}
where $T^{c,\text{iso}} = \mathbf{S}:(\D\lambda^*, \D\mu^*_{\text{off}})$
is the isotropic Eshelby contraction and
$\delta\mu^* = \D\mu^*_{\text{diag}} - \D\mu^*_{\text{off}}$.
The off-diagonal shear $T_2^c$ receives no cubic correction because
$E_{ijkl}$ vanishes for $i \neq j$.  These amplified values are
the final physical T-matrix---no additional denominator correction
is applied.

\subsection{Physical T-matrix mapping}
\label{sec:physical-mapping}

The per-irrep scalars map to the physical T-matrix coefficients via:
\begin{align}
  T_1^c &= \tfrac{1}{3}\big(\sigma_{A_{1g}} - \sigma_{E_g}\big), \label{eq:T1c-map} \\
  T_2^c &= \tfrac{1}{2}\,\sigma_{T_{2g}}, \label{eq:T2c-map} \\
  T_3^c &= \sigma_{E_g} - \sigma_{T_{2g}}. \label{eq:T3c-map}
\end{align}
The inverse relations are
\begin{align}
  \sigma_{A_{1g}} &= 3T_1^c + 2T_2^c + T_3^c, \\
  \sigma_{E_g}    &= 2T_2^c + T_3^c, \\
  \sigma_{T_{2g}} &= 2T_2^c.
\end{align}

\subsection{Born limit}
\label{sec:born-limit}

At Born level (first order in contrasts), all amplification
factors reduce to unity ($A = 1$), and the amplified
T-matrix~\eqref{eq:T1c-amp}--\eqref{eq:T3c-amp} reduces to
the bare Eshelby contraction:
\begin{equation}
  T_i^{c,\text{Born}} = T_i^{c,0}
  = \mathbf{S}:(\D\lambda, \D\mu) + C^c\text{-correction}.
  \label{eq:born-T}
\end{equation}
The density-contrast body force enters only in the ungerade sector
and does not affect $T_1^c$, $T_2^c$, $T_3^c$.
The cubic anisotropy $T_3^c$ arises from the $C^c$ component
of the Eshelby integral.

\subsection{Cubic anisotropy coefficient}

The cubic anisotropy $T_3^c = \sigma_{E_g} - \sigma_{T_{2g}}$ arises
from the $C^c$ component of the Eshelby integral, which captures the
non-spherical geometry of the cube.  At Born level:
\begin{equation}
  T_3^{c,\text{Born}}
  = (2T_2^c + T_3^c) - 2T_2^c
  = T_3^c,
\end{equation}
which is determined entirely by the cubic Eshelby integral~$C^c$.
For a sphere, $C^c = 0$ and $T_3^c$ vanishes identically.
In the Galerkin framework, this anisotropy enters through the
$B$-channel of the body bilinear form (for ungerade modes) and
through $C^c$ in the Eshelby tensor (for gerade modes).

\subsection{Displacement and quadratic sector (ungerade irreps)}
\label{sec:ungerade-sector}

The four ungerade irreps ($T_{1u}$, $T_{2u}$, $A_{2u}$, $E_u$) carry
the displacement and quadratic modes.  The $T_{1u}$ block is physically
the most important: its leading eigenvalue is the \emph{dipole mode},
which encodes the density-contrast scattering.

\paragraph{Per-irrep T-matrix formula.}
For an ungerade irrep $\rho$ with multiplicity $m_\rho$, the
$m_\rho \times m_\rho$ T-matrix block is
\begin{equation}
  T_\rho
  = \big(M_\rho + B^{\text{LS}}_\rho - \varepsilon\,B_{\text{body},\rho}\big)^{-1}
    \big(\varepsilon\,B_{\text{body},\rho} - B^{\text{LS}}_\rho\big),
  \label{eq:ungerade-T}
\end{equation}
where $\varepsilon = \omega^2 \D\rho$ is the density coupling,
$M_\rho$ is the mass matrix, $B_{\text{body},\rho}$ the body bilinear
form, and $B^{\text{LS}}_\rho$ the LS-convolved stiffness.  All three
matrices are real, symmetric, and dimensionless.  For the scalar
irreps ($m_\rho = 1$), this reduces to
\begin{equation}
  \sigma_\rho
  = \frac{\varepsilon\,b_\rho - b^{\text{LS}}_\rho}
         {M_\rho + b^{\text{LS}}_\rho - \varepsilon\,b_\rho}.
\end{equation}

\paragraph{Body bilinear form.}
The body bilinear separates into isotropic ($A$) and anisotropic ($B$)
channels via the Kelvin decomposition of the Green's tensor:
\begin{equation}
  B_{\text{body},\rho}
  = a_0\,B^{(A)}_\rho + b_0\,B^{(B)}_\rho,
  \qquad
  a_0 = \frac{\alpha^2 + \beta^2}{8\pi\rho_0\alpha^2\beta^2},\quad
  b_0 = \frac{\alpha^2 - \beta^2}{8\pi\rho_0\alpha^2\beta^2}.
  \label{eq:kelvin-ab}
\end{equation}
The matrices $B^{(A)}_\rho$ and $B^{(B)}_\rho$ are numerical constants
determined by the cube geometry, computed from the
Pell-number closed forms of the body integrals (Sections~5--8).

\paragraph{LS-convolved stiffness.}
The stiffness contribution is \emph{not} the direct strain energy
$\int\boldsymbol{\varepsilon}:\D\mathbf{c}:\boldsymbol{\varepsilon}\,dV$
(which has dimensions of~Pa), but the LS-convolved form
\begin{equation}
  B^{\text{LS}}_{\alpha\beta}
  = \big\langle\bphi_\alpha
      \,\big|\, \bG \ast
      \big[\nabla\!\cdot\!(\D\mathbf{c}:\be(\bphi_\beta))\big]
    \big\rangle,
  \label{eq:bstiff-ls}
\end{equation}
which is dimensionless.  Here $\bG$ is the static Green's tensor and
$\be(\bphi_\beta) = \tfrac{1}{2}(\nabla\bphi_\beta + \nabla\bphi_\beta^T)$
is the strain of the trial function.

For a quadratic mode $\bphi_\beta$ on the cube $\Omega = [-1,1]^3$,
the divergence of the stress perturbation splits into a volume force
and a surface traction:
\begin{equation}
  \nabla\!\cdot\!(\D\mathbf{c}:\be(\bphi_\beta))
  = \mathbf{f}_{\text{vol}}^{(\beta)}\,\chi_\Omega
  + \mathbf{t}^{(\beta)}\,\delta_{\partial\Omega},
  \label{eq:vol-surf-split}
\end{equation}
where $\chi_\Omega$ is the indicator function and
$\delta_{\partial\Omega}$ the surface delta.  The volume force is
constant inside $\Omega$:
\begin{equation}
  \mathbf{f}_{\text{vol}} = (\D\lambda + \D\mu)\,\nabla(\nabla\!\cdot\!\bphi_\beta)
  + \D\mu\,\nabla^2\bphi_\beta.
\end{equation}
For S-type quadratic modes $r_a^2\,\mathbf{e}_p$:
\begin{equation}
  \mathbf{f}_{\text{vol}} = \begin{cases}
    2(\D\lambda + 2\D\mu)\,\mathbf{e}_p & \text{if } a = p, \\
    2\D\mu\,\mathbf{e}_p                & \text{if } a \neq p,
  \end{cases}
\end{equation}
and for X-type modes $r_b r_c\,\mathbf{e}_p$ (with $b \neq c$),
$\mathbf{f}_{\text{vol}} = 0$.  The surface traction
$\mathbf{t}^{(\beta)} = (\D\mathbf{c}:\be(\bphi_\beta))\cdot\hat{\mathbf{n}}$
is at most linear in the face coordinates.

This decomposition is computationally advantageous.  The volume part of
the LS convolution reduces to existing body integrals:
\begin{equation}
  B^{\text{LS,vol}}_{\alpha\beta}
  = \sum_{j=1}^{3} f^{(\beta)}_j \, B_{\text{body}}[\alpha,\,\mathbf{e}_j],
  \label{eq:bstiff-vol}
\end{equation}
where $B_{\text{body}}[\alpha,\,\mathbf{e}_j]$ is the body bilinear
between $\bphi_\alpha$ and constant displacement $\mathbf{e}_j$---the
first three columns of the body matrix already computed in Sections~5--8.
The surface part~$B^{\text{LS,surf}}$ is a five-dimensional integral
(3D from $\mathbf{r} \in \Omega$, 2D from $\mathbf{r}'$ on a cube
face), with a $1/|\mathbf{r}-\mathbf{r}'|$ singularity that is
weaker than the 6D body bilinear.  A key result of this work
(Section~\ref{sec:analytical-reduction}) is that this integral,
like the body bilinear, reduces analytically to the same family of 3D
master integrals $\mathcal{M}^+_{p,q,r}$ and $\mathcal{M}^{B+}_{p,q,r}$
via the $(\sigma,\xi,v)$ substitution, requiring no numerical quadrature.

\paragraph{Four-channel decomposition.}
The LS-convolved stiffness is bilinear in
$(a_0, b_0) \times (\D\lambda, \D\mu)$, giving a decomposition into
four independent numerical matrices:
\begin{equation}
  B^{\text{LS}}_\rho
  = a_0 \D\lambda\,S^{(A\lambda)}_\rho
  + a_0 \D\mu\,S^{(A\mu)}_\rho
  + b_0 \D\lambda\,S^{(B\lambda)}_\rho
  + b_0 \D\mu\,S^{(B\mu)}_\rho.
  \label{eq:four-channel}
\end{equation}
These four matrices per irrep are geometry-dependent constants of the
cube, analogous to the $A$- and $B$-channel body matrices.  They are
computed by \texttt{CubeT27Stiffness\_LS.wl}.

\paragraph{Physical content of the $T_{1u}$ block.}
The $T_{1u}$ irrep appears with multiplicity $m = 4$: one constant
displacement mode and three quadratic modes (two S-type, one X-type).
Its leading eigenvalue is the \emph{dipole mode}, which at low
frequency ($ka \ll 1$) dominates all other ungerade eigenvalues
by a factor $\gtrsim 3$.  This mode encodes the rigid-body response
of the scatterer to density contrast and is the elastic analogue of
the acoustic monopole.  The three subdominant eigenvalues are
quadratic corrections that become significant only at $ka \gtrsim 0.3$.

The $T_{2u}$ ($2 \times 2$) block carries two pure quadratic modes,
while $A_{2u}$ and $E_u$ are $1 \times 1$ pure quadratic scalars.
For a sphere, all quadratic eigenvalues would be degenerate within
each parity class; the cubic geometry of the cube lifts this degeneracy
through the $B$-channel.

\subsection{Unified analytical reduction: all integrals reduce to
master integrals}
\label{sec:analytical-reduction}

The central computational result of this work is that \emph{every}
entry of the $27 \times 27$ Galerkin closure---both the six-dimensional
body bilinear forms and the five-dimensional surface stiffness
integrals---reduces analytically to finite linear combinations of two
families of three-dimensional master integrals over the unit cube
$[0,1]^3$:
\begin{align}
  \mathcal{M}^+_{p,q,r}
    &:= \int_{[0,1]^3}\!\frac{x^p\,y^q\,z^r}
         {\sqrt{x^2+y^2+z^2}}\,dx\,dy\,dz
    &&\text{($A$-channel)},
    \label{eq:Mp-def} \\
  \mathcal{M}^{B+}_{p,q,r}
    &:= \int_{[0,1]^3}\!\frac{x^p\,y^q\,z^r}
         {(x^2+y^2+z^2)^{3/2}}\,dx\,dy\,dz
    &&\text{($B$-channel)}.
    \label{eq:MpB-def}
\end{align}
No numerical quadrature is required at any stage: the entire Galerkin
matrix is determined by exact arithmetic once the master integrals are
evaluated (Tables~\ref{tab:tier1-masters}--\ref{tab:tier2-masters}).
Mathematica's symbolic \texttt{Integrate} returns closed forms for all
needed instances in terms of algebraic numbers, $\pi$, logarithms, and
inverse trigonometric functions.

\subsubsection*{Body bilinear form}

The reduction of the body bilinear form to master integrals was
described in
Sections~\ref{sec:cm-relative}--\ref{sec:results-tier3}.  In summary:
the centre-of-mass/relative coordinate substitution
$(\bm{\sigma},\bm{s})$ (eq.~\eqref{eq:cm-relative}) collapses the
six-dimensional double integral to a three-dimensional integral against
a polynomial weight.  After the octant-symmetry restriction
$\bm{u} \in [0,1]^3$, every entry of
$\mathcal{B}_{\mathrm{body}}$ is a signed linear combination of
$\mathcal{M}^+_{p,q,r}$ values (for the isotropic $A$-channel,
kernel $\delta_{ij}/|\bm{u}|$) and of $B$-channel integrals involving
$u_i u_j / |\bm{u}|^3$ (reducible to $\mathcal{M}^{B+}_{p,q,r}$).
The maximum monomial degree is $p+q+r \le 3$ for the $9{\times}9$
linear sub-block and $p+q+r \le 7$ for the full $27{\times}27$ body
matrix.  This reduction extends the classical Eshelby
integrals~\cite{mura1987} from a single polynomial weight (the cube
autocorrelation) to arbitrary polynomial trial-function pairs, and is
the basis for the closed-form results in
Tables~\ref{tab:tier1-masters}--\ref{tab:high-prec-scalars}.

\subsubsection*{Surface stiffness integrals: the
$(\sigma,\xi,v)$ substitution}

The surface part of the LS-convolved stiffness,
\begin{equation}
  B^{\text{LS,surf}}_{\alpha\beta}
  = \int_V\!\!\oint_{\partial V}\!\varphi_{\alpha,i}(\br)\,
    G_{in}(\br-\br')\,t^{(\beta)}_n(\br')\,dS'\,dV,
  \label{eq:bstiff-surf-def}
\end{equation}
involves a five-dimensional integral: three from $\br \in V = [-1,1]^3$
and two from $\br'$ on a face of~$V$.  The crucial observation is that
this integral admits the \emph{same} centre-of-mass/relative change of
variables used for the body bilinear, supplemented by a normal-distance
coordinate.

Treating each face separately, the integral on the face
$r_k = +1$ (the opposite face follows by reflection) uses
\begin{equation}
\begin{aligned}
  \sigma_a &= \tfrac{1}{2}(r_a + s_a), &
  \xi_a &= \tfrac{1}{2}(r_a - s_a), &
  &\text{for each shared axis } a \neq k,\\
  v &= \tfrac{1}{2}(1 - r_k), & & &
  &\text{(distance from face, $v \in [0,1]$).}
\end{aligned}
\label{eq:sigma-xi-v}
\end{equation}
Here $s_a$ denotes the face coordinates of $\br'$, and the field point
has $r_k = 1 - 2v$.  In these coordinates the separation vector and
its norm become
\begin{equation}
  \br - \br'
  = \bigl(\ldots,\,2\xi_a,\,\ldots,\,-2v,\,\ldots\bigr),
  \qquad
  |\br-\br'| = 2\rho,
  \qquad
  \rho := \sqrt{v^2 + \xi_1^2 + \xi_2^2},
  \label{eq:face-separation}
\end{equation}
where the $2\xi_a$ entries occupy the shared-axis positions and $-2v$
the normal-axis position.  The Jacobian of the substitution
\eqref{eq:sigma-xi-v} is a constant ($= 8$), and the integrand of
\eqref{eq:bstiff-surf-def} factorises as
\begin{equation}
  \underbrace{P(\sigma, \xi, v)}_{\text{polynomial}}
  \;\times\;
  \begin{cases}
    \dfrac{1}{2\rho} &\text{($A$-channel)},\\[6pt]
    \dfrac{x_i\,x_j}{(2\rho)^3} &\text{($B$-channel)}.
  \end{cases}
  \label{eq:face-factorisation}
\end{equation}
Because the polynomial part is at most quadratic in $\sigma_a$ and the
integration domain in $\sigma_a$ is the symmetric interval
$[-(1-|\xi_a|),\,1-|\xi_a|]$, the $\sigma$-integration is elementary:
odd powers of $\sigma_a$ vanish by parity, and even powers yield
\begin{equation}
  \int_{-(1-|\xi_a|)}^{1-|\xi_a|}\! \sigma_a^{2j}\,d\sigma_a
  \;=\; \frac{2\,(1-|\xi_a|)^{2j+1}}{2j+1}.
  \label{eq:sigma-int}
\end{equation}
After this step, the $\sigma$-variables are eliminated and the integral
over each face reduces to a three-dimensional integral over
$(\xi_1,\xi_2,v) \in [0,1]^3$ (using octant symmetry
$|\xi_a| \to \xi_a$).  Expanding the polynomial in monomials
$\xi_1^p\,\xi_2^q\,v^r$ and reading off the coefficients gives,
for each face and channel,
\begin{equation}
  B^{\text{LS,surf}}_{\alpha\beta}\big|_{\text{face}}
  \;=\; 64\!\!\sum_{p,q,r}\! c^{(A)}_{pqr}\,\mathcal{M}^+_{p,q,r}
  \;+\; 16\!\!\sum_{p,q,r}\!c^{(B)}_{pqr}\,\mathcal{M}^{B+}_{p,q,r},
  \label{eq:surf-master-expansion}
\end{equation}
where the coefficients $c^{(A,B)}_{pqr}$ are rational numbers
determined by the mode polynomials $\bphi_\alpha$, the traction
$\mathbf{t}^{(\beta)}$, and the face orientation.
The prefactor $64 = 8 \times 4 \times 2$ absorbs the Jacobian ($8$),
the octant symmetry in $(\xi_1,\xi_2)$ ($4$), and the two paired faces
per axis ($2$); the $B$-channel prefactor $16$ reflects the additional
$1/(2\rho)^2$ relative to the $A$ kernel.  The maximum monomial degree
required is $p+q+r \le 4$ for the surface integrals on $T_{27}$.
Summing over the three coordinate axes gives the full surface
stiffness.

\subsubsection*{Completeness of the reduction}

Combining the body and surface results, the full LS-convolved
stiffness \eqref{eq:bstiff-ls} on the $T_{27}$ trial space is
\begin{equation}
  B^{\text{LS}}_{\alpha\beta}
  = B^{\text{LS,vol}}_{\alpha\beta}
  + B^{\text{LS,surf}}_{\alpha\beta}
  = \text{(linear combination of $\mathcal{M}^+$ and
    $\mathcal{M}^{B+}$ values)}.
  \label{eq:bstiff-full-analytical}
\end{equation}
Together with the body bilinear form and the mass matrix (which is a
polynomial integral and hence trivially exact), the \emph{entire}
$27 \times 27$ Galerkin closure is determined by the master integrals
$\mathcal{M}^+_{p,q,r}$ and $\mathcal{M}^{B+}_{p,q,r}$.  The two
families are moreover linked by the trace identity
\begin{equation}
  \mathcal{M}^{B+}_{p+2,q,r} + \mathcal{M}^{B+}_{p,q+2,r}
  + \mathcal{M}^{B+}_{p,q,r+2}
  = \mathcal{M}^+_{p,q,r},
  \label{eq:MpB-trace-identity}
\end{equation}
which follows from $|x|^2 / |x|^3 = 1/|x|$ and reduces the number
of independent $B$-channel evaluations.  The total number of distinct
master integrals needed is $\sim\!30$ (after cubic-symmetry
permutations and the trace identity); Mathematica's
\texttt{Integrate} evaluates each one in closed form in $5$--$30$ s,
with a total computation time of $\sim\!12$ minutes for the complete
analytical evaluation of all $27 \times 27$ stiffness entries.  This
is a $\sim\!300{\times}$ speedup over the previous $5$D
\texttt{NIntegrate} approach ($\sim\!60$ min), and---critically---the
analytical values are exact rather than approximate.

\begin{remark}[Catastrophic failure of direct numerical quadrature]
\label{rem:ninteg-failure}
{\sloppy
The analytical reduction was motivated by the discovery that
Mathematica's five-dimensional \texttt{NIntegrate} with
\texttt{DuffyCoordinates} singularity handling gave
\emph{catastrophically wrong} results for the surface stiffness
integrals involving quadratic mode pairs.  The error mechanism is
severe cancellation: after the $(\sigma,\xi,v)$ substitution, the
$\sigma$-integrated polynomial\par}
\begin{equation*}
  \operatorname{sigIntA}(\xi_1,\xi_2,v)
  = \tfrac{4}{3}(1-2v)(1-\xi_1)(1-\xi_2)
    (1-\xi_1-\xi_2-\xi_1^2-\xi_2^2)
\end{equation*}
(shown here for the $E_u$ entry, $A$-channel) is the product of a
tent-like envelope and a polynomial that changes sign.  The resulting
integral is $\sim\!3.02$, but the individual monomial contributions
involve $\mathcal{M}^+$ values of order $\sim\!100$ with alternating
signs.  Adaptive quadrature in five dimensions resolves the singularity
but not the cancellation: the \texttt{NIntegrate} result was $136.0$,
a factor-of-$45$ error.  The linear--linear sub-block had smaller
errors ($0.2$--$3\%$), and the linear--quadratic cross terms $1$--$3\%$.

Independent verification of the analytical values was carried out by
three methods:
\begin{enumerate}[label=(\roman*)]
  \item Monte Carlo integration ($5 \times 10^7$ samples) on the
    original 5D integral, giving $3.017 \pm 0.007$;
  \item \texttt{scipy.integrate.tplquad} on the reduced 3D integral,
    giving $64 \times 0.04715 = 3.018$;
  \item direct summation $64\sum c_{pqr}\,\mathcal{M}^+_{p,q,r}$ using
    independently computed master integrals, giving $3.018$.
\end{enumerate}
All three agree; the \texttt{NIntegrate} result of $136.0$ is
definitively incorrect.
\end{remark}

%======================================================================
\subsection{Plane-wave scattering from the \texorpdfstring{$T_{27}$}{T27}
closure}
\label{sec:plane-wave}
%======================================================================

With all $27 \times 27$ Galerkin matrices determined analytically, we now
derive the complete plane-wave scattering framework for a single cube in
a homogeneous elastic background.  The result is a frequency-dependent
$27 \times 27$ T-matrix that maps an incident plane-wave field to scattered
multipole sources, with exact P--SV--SH decomposition in the far field.

\subsubsection{Assembly of the full
\texorpdfstring{$27 \times 27$}{27x27} T-matrix}
\label{sec:T27-assembly}

The symmetry-adapted basis change $U_{\mathrm{sym}}$ of Section~9c
block-diagonalises all three Galerkin matrices ($M_{27}$,
$B_{\mathrm{body}}$, $B_{\mathrm{el}}$) simultaneously.  The full
T-matrix in the physical (polynomial) basis is
\begin{equation}
  \boxed{\;T_{27}
  \;=\; U_{\mathrm{sym}}^{\top}\;
  \mathcal{T}_{\mathrm{diag}}\;
  U_{\mathrm{sym}}\;}
  \label{eq:T27-assembly}
\end{equation}
where $\mathcal{T}_{\mathrm{diag}}$ is the $27 \times 27$
block-diagonal matrix in the irrep-adapted basis.  The parity
decomposition $T_{27} = T_9 \oplus T_{27}^{\mathrm{quad}}$
(constant+linear $\oplus$ quadratic) splits further by O$_h$ irrep
into an ungerade sector (21 dimensions: displacement $+$ quadratic
modes) and a gerade sector (6 dimensions: strain modes):
\begin{equation}
  \mathcal{T}_{\mathrm{diag}}
  \;=\;
  \begin{pmatrix}
    \mathcal{T}_u & 0 \\[2pt]
    0 & \mathcal{T}_g
  \end{pmatrix}_{\!27\times 27}.
  \label{eq:T-parity-split}
\end{equation}

\paragraph{Ungerade sector ($21 \times 21$).}
The ungerade block decomposes into four irreps.  By Schur's lemma,
each irrep~$\rho$ of dimension~$d_\rho$ and multiplicity~$m_\rho$
contributes $d_\rho$ \emph{identical} copies of the $m_\rho \times
m_\rho$ per-irrep T-matrix:
\begin{equation}
  \mathcal{T}_u
  \;=\;
  \begin{pmatrix}
    \mathcal{T}_{T_{1u}} & & & \\
    & \mathcal{T}_{T_{2u}} & & \\
    & & \sigma_{A_{2u}} & \\
    & & & \mathcal{T}_{E_u}
  \end{pmatrix}_{\!21\times 21},
  \label{eq:Tu-block}
\end{equation}
where the sub-blocks are
\begin{equation}
  \mathcal{T}_{T_{1u}}
  =
  \begin{pmatrix}
    T^{(T_{1u})} & & \\
    & T^{(T_{1u})} & \\
    & & T^{(T_{1u})}
  \end{pmatrix}_{\!12\times 12}
  \!\!\!,\qquad
  \mathcal{T}_{T_{2u}}
  =
  \begin{pmatrix}
    T^{(T_{2u})} & & \\
    & T^{(T_{2u})} & \\
    & & T^{(T_{2u})}
  \end{pmatrix}_{\!6\times 6}
  \!\!\!,
  \label{eq:Tu-subblocks}
\end{equation}
\begin{equation}
  \mathcal{T}_{E_u}
  = \sigma_{E_u}\,I_2\,.
  \label{eq:TEu-scalar}
\end{equation}
The three copies within each triplet ($T_{1u}$, $T_{2u}$) correspond
to the three Cartesian components of the irreducible triplet; the two
copies in $E_u$ correspond to the doublet.  The single $A_{2u}$ scalar
$\sigma_{A_{2u}}$ stands alone.

The per-irrep T-matrix $T^{(\rho)}$ for each ungerade irrep is an
$m_\rho \times m_\rho$ matrix given by
\begin{equation}
  \boxed{\;T^{(\rho)}
  \;=\; \bigl(M_\rho + B^{\mathrm{LS}}_\rho
    - \varepsilon\,B_{\mathrm{body},\rho}\bigr)^{-1}
    \bigl(\varepsilon\,B_{\mathrm{body},\rho}
    - B^{\mathrm{LS}}_\rho\bigr),\;}
  \label{eq:T-ungerade}
\end{equation}
where $\varepsilon = \omega^2\D\rho$ is the inertial coupling and the
three matrices appearing in \eqref{eq:T-ungerade} are:
\begin{itemize}
\item $M_\rho$: the $m_\rho \times m_\rho$ mass matrix (rational,
  exact);
\item $B_{\mathrm{body},\rho} = a_0\,B^{(A)}_\rho + b_0\,B^{(B)}_\rho$:
  the body bilinear form, with Kelvin coefficients $a_0$, $b_0$
  from~\eqref{eq:kelvin-ab} and numerical matrices $B^{(A,B)}_\rho$
  determined by the master integrals of
  Section~\ref{sec:analytical-reduction};
\item $B^{\mathrm{LS}}_\rho$: the LS-convolved stiffness, decomposing
  into four channels (eq.~\eqref{eq:four-channel}):
  \begin{equation}
    B^{\mathrm{LS}}_\rho
    = a_0\D\lambda\,S^{(A\lambda)}_\rho
    + a_0\D\mu\,S^{(A\mu)}_\rho
    + b_0\D\lambda\,S^{(B\lambda)}_\rho
    + b_0\D\mu\,S^{(B\mu)}_\rho.
    \label{eq:BLS-four-channel}
  \end{equation}
\end{itemize}
For $T_{1u}$ ($m = 4$), the matrices $M$, $B^{(A,B)}$, and
$S^{(A\lambda,A\mu,B\lambda,B\mu)}$ are each $4 \times 4$, all
computed analytically from the master integrals.  For $T_{2u}$
($m = 2$), they are $2 \times 2$.  For $A_{2u}$ and $E_u$ ($m = 1$),
the formula~\eqref{eq:T-ungerade} reduces to the scalar
\begin{equation}
  \sigma_\rho
  = \frac{\varepsilon\,b_\rho - b^{\mathrm{LS}}_\rho}
         {M_\rho + b^{\mathrm{LS}}_\rho - \varepsilon\,b_\rho}\,,
  \qquad \rho \in \{A_{2u},\,E_u\}.
  \label{eq:T-ungerade-scalar}
\end{equation}

\paragraph{Gerade sector ($6 \times 6$).}
The three gerade irreps are all $m_\rho = 1$ (scalar blocks),
determined by the self-consistent Eshelby amplification of
eqs.~\eqref{eq:T1c-amp}--\eqref{eq:T3c-amp}:
\begin{equation}
  \mathcal{T}_g
  \;=\;
  \begin{pmatrix}
    \sigma_{A_{1g}} & & \\[2pt]
    & \sigma_{E_g}\,I_2 & \\[2pt]
    & & \sigma_{T_{2g}}\,I_3
  \end{pmatrix}_{\!6\times 6},
  \label{eq:Tg-block}
\end{equation}
where $I_n$ is the $n \times n$ identity and the per-irrep
scalars are
\begin{equation}
  \boxed{\;\sigma_\rho
  = \frac{\lambda_\rho(\D\mathbf{c}^*)}
         {1 - \lambda_\rho(\D\mathbf{c}^*)},
  \qquad
  \rho \in \{A_{1g},\,E_g,\,T_{2g}\},\;}
  \label{eq:T-gerade}
\end{equation}
with $\lambda_\rho$ the per-irrep eigenvalue of the amplified
Eshelby contraction and $\D\mathbf{c}^*$ the self-consistent
stiffness contrast.  Explicitly (from
eqs.~\eqref{eq:born-sigma-irrep}--\eqref{eq:Dlambda-star}):
\begin{align}
  \lambda_{A_{1g}} &= 3T_1^c + 2T_2^c + T_3^c
    &&\text{(volumetric)}, \label{eq:lam-A1g}\\
  \lambda_{E_g}    &= 2T_2^c + T_3^c
    &&\text{(deviatoric diagonal)}, \label{eq:lam-Eg}\\
  \lambda_{T_{2g}} &= 2T_2^c
    &&\text{(off-diagonal shear)}, \label{eq:lam-T2g}
\end{align}
where $T_1^c$, $T_2^c$, $T_3^c$ are the amplified coupling
coefficients.  The cubic anisotropy of the cube enters through
$T_3^c \neq 0$, splitting $\sigma_{E_g}$ from $\sigma_{T_{2g}}$
(cf.\ eq.~\eqref{eq:T3c-signature}).

\paragraph{Dimension count.}
\begin{equation}
  \underbrace{3 \times 4}_{T_{1u}} \;+\;
  \underbrace{3 \times 2}_{T_{2u}} \;+\;
  \underbrace{1}_{A_{2u}} \;+\;
  \underbrace{2}_{E_u} \;+\;
  \underbrace{1}_{A_{1g}} \;+\;
  \underbrace{2}_{E_g} \;+\;
  \underbrace{3}_{T_{2g}}
  \;=\; 12 + 6 + 1 + 2 + 1 + 2 + 3
  \;=\; 27.
  \label{eq:dim-count}
\end{equation}

\paragraph{Consistency check.}
The leading $9 \times 9$ sub-block of $T_{27}$ (rows/columns~$1$--$9$,
corresponding to the constant and linear basis functions of~$T_9$)
is the Schur-complemented T-matrix on~$T_9$: it incorporates the
effect of the quadratic modes through the body-channel Schur
complement of Section~12b.  This sub-block must agree with the Path-A
T-matrix from \texttt{compute\_cube\_tmatrix()} to the Galerkin
truncation error ($\lesssim 1\%$ at $ka = 0.3$).

\subsubsection{Incident plane-wave projection}
\label{sec:incident-projection}

An elastic plane wave in the background medium has the form
\begin{equation}
  \bu^0(\br) = \hat{\bm{d}}\,\exp\!\bigl(i\,\bm{k}\cdot\br\bigr),
  \qquad
  \bm{k} = \frac{\omega}{c}\,\hat{\bm{k}},
  \label{eq:plane-wave}
\end{equation}
where $\hat{\bm{d}}$ is the polarisation, $\hat{\bm{k}}$ the
propagation direction, and $c = \alpha$ for a P-wave ($\hat{\bm{d}} = \hat{\bm{k}}$),
$c = \beta$ for an S-wave ($\hat{\bm{d}} \perp \hat{\bm{k}}$).
Projecting onto the Galerkin basis $\{\bphi_\alpha\}$ of $T_{27}$
centred at the origin gives the incident-field coefficient vector
\begin{equation}
  c^0_\alpha
  = \bigl[M_{27}^{-1}\bigr]_{\alpha\beta}\!
    \int_{V}\!\bphi_\beta(\br)\cdot\bu^0(\br)\,dV.
  \label{eq:c0-projection}
\end{equation}
Because every $\bphi_\beta$ is a monomial $r_1^{n_1}r_2^{n_2}r_3^{n_3}\be_i$,
the overlap integral factorises over the three axes of the cube
$V = [-a,a]^3$:
\begin{equation}
  \int_{V}\!r_1^{n_1}r_2^{n_2}r_3^{n_3}\,e^{i\bm{k}\cdot\br}\,dV
  = \prod_{j=1}^{3}\mathcal{I}_{n_j}(k_j,a),
  \label{eq:overlap-factorise}
\end{equation}
where the one-dimensional overlap integrals satisfy the recursion
\begin{equation}
  \mathcal{I}_{n}(k,a)
  := \int_{-a}^{a}\!x^n\,e^{ikx}\,dx
  = \frac{a^n\bigl(e^{ika} - (-1)^n e^{-ika}\bigr)}{ik}
    - \frac{n}{ik}\,\mathcal{I}_{n-1}(k,a),
  \label{eq:overlap-recursion}
\end{equation}
with $\mathcal{I}_{0}(k,a) = 2a\operatorname{sinc}(ka)$.  The first
few are
\begin{align}
  \mathcal{I}_0 &= 2a\operatorname{sinc}(ka), \label{eq:I0}\\
  \mathcal{I}_1 &= \frac{2i\bigl[\sin(ka) - ka\cos(ka)\bigr]}{k^2},
  \label{eq:I1}\\
  \mathcal{I}_2 &= \frac{2\bigl[(2-k^2a^2)\sin(ka) - 2ka\cos(ka)\bigr]}{k^3}.
  \label{eq:I2}
\end{align}

\paragraph{Low-frequency expansion.}
For $ka \ll 1$ the overlaps simplify to the polynomial moments of the
cube:
\begin{align}
  \mathcal{I}_0 &\approx 2a, &
  \mathcal{I}_1 &\approx \tfrac{2i}{3}\,k a^3, &
  \mathcal{I}_2 &\approx \tfrac{2}{3}\,a^3.
  \label{eq:overlap-lowfreq}
\end{align}
The incident-field vector then has the physical interpretation:
\begin{equation}
  c^0_\alpha \;\xrightarrow{ka\to 0}\;
  \begin{cases}
    d_i                               & \alpha = i \in \{1,2,3\}
      \quad\text{(displacement at origin)},\\[4pt]
    \varepsilon^0_{ij}(\mathbf{0})    & \alpha \leftrightarrow (i,j) \in \{4,\ldots,9\}
      \quad\text{(incident strain)},\\[4pt]
    \tfrac{1}{2}\partial_a\partial_b\,u^0_i(\mathbf{0})
                                      & \alpha \in \{10,\ldots,27\}
      \quad\text{(strain gradient)},
  \end{cases}
  \label{eq:c0-lowfreq}
\end{equation}
so that the constant modes carry the displacement monopole at
$O(1)$, the linear modes the strain at $O(ka)$, and the
quadratic modes the strain gradient at $O((ka)^2)$.  This
hierarchy is the basis for the Born-order multipole expansion
of the scattered field.

\paragraph{The three wave types.}
For a propagation direction $\hat{\bm{k}}$, the three independent
polarisations are:
\begin{align}
  \text{P:}\quad  &\hat{\bm{d}} = \hat{\bm{k}},
    &&c = \alpha, \label{eq:P-wave}\\
  \text{SV:}\quad &\hat{\bm{d}} = \hat{\bm{\theta}}
    \;\;\text{(in the $\hat{\bm{k}}$--$\hat{\bm{z}}$ plane,
    $\perp\hat{\bm{k}}$)},
    &&c = \beta, \label{eq:SV-wave}\\
  \text{SH:}\quad &\hat{\bm{d}} = \hat{\bm{\phi}}
    \;\;\text{(perpendicular to the $\hat{\bm{k}}$--$\hat{\bm{z}}$ plane)},
    &&c = \beta. \label{eq:SH-wave}
\end{align}
In the coordinate system of this paper ($z = $ axis~0, down;
$x = $ axis~1, right; $y = $ axis~2, out), a wave propagating
vertically downward has $\hat{\bm{k}} = \hat{\bm{z}}$,
$\hat{\bm{\theta}} = \hat{\bm{x}}$,
$\hat{\bm{\phi}} = \hat{\bm{y}}$.

\subsubsection{Scattered source: multipole decomposition}
\label{sec:multipole}

The T-matrix maps the incident-field vector to scattered-source
coefficients:
\begin{equation}
  \bm{c}^{\mathrm{sc}} = T_{27}\,\bm{c}^0.
  \label{eq:csc-def}
\end{equation}
The 27 scattered coefficients decompose into three physical multipole
orders, each of which radiates with a distinct angular pattern:
\begin{equation}
\renewcommand{\arraystretch}{1.4}
\begin{array}{lcll}
\toprule
\text{Coefficients} & \text{Dim} & \text{Multipole source}
  & \text{Radiation order} \\
\midrule
c^{\mathrm{sc}}_{1\text{--}3}
  & 3 & \text{Force monopole } F_i
  & \text{Dipole, } \ell = 1 \\
c^{\mathrm{sc}}_{4\text{--}9}
  & 6 & \text{Stress dipole } D_{ij}
  & \text{Quadrupole, } \ell = 2 \\
c^{\mathrm{sc}}_{10\text{--}27}
  & 18 & \text{Force quadrupole } Q_{ijk}
  & \text{Octupole, } \ell = 3 \\
\bottomrule
\end{array}
\label{eq:multipole-table}
\end{equation}

The physical source amplitudes are obtained by contracting the scattered
coefficients with the material contrasts and cube volume $V = (2a)^3$:
\begin{align}
  F_i &= \omega^2\D\rho\,V\,c^{\mathrm{sc}}_i
    + \D c_{ijkl}\!\oint_{\partial V}\!
    \varepsilon_{kl}(\bu^{\mathrm{sc}})\,n_j\,dS,
    \label{eq:force-monopole}\\[3pt]
  D_{ij} &= \int_V\!\bigl[
    \omega^2\D\rho\,u^{\mathrm{sc}}_i\,r_j
    + \sigma^{\mathrm{sc}}_{ij}\bigr]\,dV,
    \label{eq:stress-dipole}
\end{align}
where $\sigma^{\mathrm{sc}}_{ij} = \D c_{ijkl}\,\varepsilon_{kl}(\bu^{\mathrm{sc}})$
is the stress perturbation.  For the constant and linear modes of $T_9$,
these reduce to the standard force--stress-dipole pair used in the
existing Foldy--Lax framework~\cite{mura1987}.  The quadratic modes
generate a force-quadrupole correction $Q_{ijk}$ that contributes at
$O((ka)^3)$.

\subsubsection{Far-field scattering amplitudes}
\label{sec:far-field}

In the radiation zone ($|\br| \gg a$), the scattered displacement
from the multipole sources separates into P and S contributions.
Using the asymptotic expansion of the Green's tensor,
\begin{equation}
  G_{ij}(\br) \;\xrightarrow{r\to\infty}\;
  G^P_{ij}(\hat{\br})\,\frac{e^{ik_P r}}{r}
  \;+\; G^S_{ij}(\hat{\br})\,\frac{e^{ik_S r}}{r},
  \label{eq:G-farfield}
\end{equation}
where $G^P_{ij} = \hat{r}_i\hat{r}_j/(4\pi\rho\alpha^2)$ and
$G^S_{ij} = (\delta_{ij} - \hat{r}_i\hat{r}_j)/(4\pi\rho\beta^2)$
are the P- and S-wave far-field projection tensors, the scattered
P-wave amplitude in the direction $\hat{\br}$ is
\begin{equation}
  f^P_i(\hat{\br})
  = -\frac{\hat{r}_i}{4\pi\rho\alpha^2}\,
    \Bigl[\hat{r}_j F_j
    \;-\; ik_P\,\hat{r}_j\hat{r}_k\,D_{jk}
    \;-\; \tfrac{1}{2}k_P^2\,\hat{r}_j\hat{r}_k\hat{r}_l\,Q_{jkl}
    \;+\; \cdots
    \Bigr],
  \label{eq:fP}
\end{equation}
and the S-wave amplitude is
\begin{equation}
  f^S_i(\hat{\br})
  = -\frac{\delta_{ij} - \hat{r}_i\hat{r}_j}{4\pi\rho\beta^2}\,
    \Bigl[F_j
    \;-\; ik_S\,\hat{r}_k\,D_{jk}
    \;-\; \tfrac{1}{2}k_S^2\,\hat{r}_k\hat{r}_l\,Q_{jkl}
    \;+\; \cdots
    \Bigr].
  \label{eq:fS}
\end{equation}
The S-wave amplitude decomposes into SV and SH components via
\begin{align}
  f^{SV}(\hat{\br}) &= \hat{\bm{\theta}} \cdot \bm{f}^S(\hat{\br}),
  \label{eq:fSV}\\
  f^{SH}(\hat{\br}) &= \hat{\bm{\phi}} \cdot \bm{f}^S(\hat{\br}),
  \label{eq:fSH}
\end{align}
where $(\hat{\br},\hat{\bm{\theta}},\hat{\bm{\phi}})$ form the
spherical coordinate triad.

\paragraph{Scattering cross sections.}
The differential scattering cross section per unit solid angle is
\begin{equation}
  \frac{d\sigma}{d\Omega}\bigg|_{\mathrm{P}\to X}
  = \frac{c_X}{c_{\mathrm{inc}}}\,
    \bigl|f^X(\hat{\br})\bigr|^2,
  \qquad X \in \{\mathrm{P},\,\mathrm{SV},\,\mathrm{SH}\},
  \label{eq:dsigma}
\end{equation}
where $c_{\mathrm{inc}}$ and $c_X$ are the incident and scattered wave
speeds, and $f^X$ is the scalar amplitude in channel~$X$.
The total scattering cross section is obtained by integrating
\eqref{eq:dsigma} over all directions, and the extinction cross section
follows from the optical theorem:
\begin{equation}
  \sigma_{\mathrm{ext}}
  = -\frac{4\pi}{k_{\mathrm{inc}}}\,
    \operatorname{Im}\!\bigl[f(\hat{\bm{k}})\bigr],
  \label{eq:optical-theorem}
\end{equation}
where $f(\hat{\bm{k}})$ is the forward scattering amplitude in the
incident-wave channel.  At low frequency the three multipole orders
contribute with distinct frequency scalings:
\begin{center}
\renewcommand{\arraystretch}{1.3}
\begin{tabular}{lll}
Force monopole (dipole): & $\sigma_{\mathrm{sc}} \sim k^4 a^6$
  & (Rayleigh) \\
Stress dipole (quadrupole): & $\sigma_{\mathrm{sc}} \sim k^4 a^6$
  & (stiffness contrast) \\
Force quadrupole (octupole): & $\sigma_{\mathrm{sc}} \sim k^6 a^8$
  & (new from $T_{27}$)
\end{tabular}
\end{center}
The dipole and quadrupole terms both scale as $k^4 a^6$ (the classic
Rayleigh scattering law) but with different angular patterns and
different sensitivity to $\D\rho$ vs.\ $\D\mu$; the octupole term from
the quadratic modes is two orders higher in $ka$, becoming significant
only in the transition regime $ka \gtrsim 0.3$.

\paragraph{Connection to the Mie limit.}
For a sphere of equivalent volume, the T-matrix partial-wave
coefficients $(a_0, a_1, a_2, \ldots)$ correspond to multipole
orders $(\ell = 0, 1, 2, \ldots)$.  The $T_{27}$ cube closure captures:
\begin{itemize}
  \item $a_0$ (monopole, $\ell = 0$): from the $A_{1g}$ irrep
    (volumetric strain),
  \item $a_1$ (dipole, $\ell = 1$): from the $T_{1u}$ irrep
    (displacement),
  \item $a_2$ (quadrupole, $\ell = 2$): from the $E_g + T_{2g}$
    irreps (deviatoric strain), split into $a_{2,E}$ and $a_{2,T}$ by
    cubic anisotropy,
  \item $a_3$ (octupole, $\ell = 3$): partially, from the quadratic
    modes ($T_{2u}$, $A_{2u}$, $E_u$, and subdominant $T_{1u}$
    eigenvalues).
\end{itemize}
In the sphere limit ($C^c \to 0$), the $E_g$ and $T_{2g}$ eigenvalues
coalesce and reproduce the isotropic $a_2$; the quadratic-mode
eigenvalues coalesce into the degenerate $\ell = 3$ channel.  The cube's
cubic anisotropy lifts these degeneracies, generating mode conversions
(P$\to$SV, SV$\to$SH) that are absent for a sphere.

\subsubsection{Validation targets for Phase~1}
\label{sec:validation-targets}

\begin{enumerate}[label=(V\arabic*)]
\item \textbf{Sphere limit.}  For the standard test parameters
  ($\alpha = 5000$, $\beta = 3000$, $\rho = 2500$; moderate contrast)
  at $ka = 0.05$, the $T_{27}$ far-field
  amplitudes must reproduce the Mie coefficients $a_0$, $a_1$, $a_2$ to
  $\sim\!1\%$ (the $O((ka)^2)$ Galerkin truncation error).
\item \textbf{Optical theorem.}  The extinction and total scattering
  cross sections computed from \eqref{eq:optical-theorem} and the
  angular integral of \eqref{eq:dsigma} must agree to $\sim\!10^{-6}$
  (numerical integration precision).
\item \textbf{$T_9$ consistency.}  The leading $9 \times 9$ sub-block
  of $T_{27}$ must reproduce the Path-A T-matrix from
  \texttt{compute\_cube\_tmatrix()} to the Galerkin truncation error
  ($\lesssim 1\%$ at $ka = 0.3$, $\lesssim 10\%$ at $ka = 0.5$).
\item \textbf{P--SV--SH power balance.}  For an incident P-wave, the
  scattered power $\sigma^{\mathrm{P\to P}} + \sigma^{\mathrm{P\to SV}}
  + \sigma^{\mathrm{P\to SH}} = \sigma_{\mathrm{ext}}^P$ must hold.
\item \textbf{Cubic symmetry.}  For incidence along $[100]$, $[010]$,
  $[001]$, the three total cross sections must be identical.  For $[110]$
  incidence, the SH cross section must vanish by mirror symmetry.
\end{enumerate}


\bibliographystyle{plain}
\begin{thebibliography}{9}
\bibitem{mura1987} T.\ Mura, \emph{Micromechanics of Defects in Solids},
   2nd ed., Martinus Nijhoff (1987).
\bibitem{khachaturyan} A.\ G.\ Khachaturyan, \emph{Theory of Structural
   Transformations in Solids}, Wiley (1983).
\bibitem{costabel-stephan} M.\ Costabel and E.\ Stephan, ``Boundary
   integral equations for mixed boundary value problems in polygonal
   domains and Galerkin approximation'', \emph{Banach Center Publ.}
   \textbf{15} (1985), 175--251.
\bibitem{hsiao-wendland} G.\ C.\ Hsiao and W.\ L.\ Wendland,
   \emph{Boundary Integral Equations}, Springer (2008).
\end{thebibliography}

\end{document}
